\section{The covolume bound}\label{sec:pf:covol}

\providecommand{\covolLcv}{\operatorname{lcv}}
\providecommand{\covolMainT}{\mathrm{mainT}}
\providecommand{\covolUnitE}{\mathrm{unitE}}
\providecommand{\covolHch}{h_{\mathrm{Ch}}}
\providecommand{\covolHform}{h_{\mathrm{F}}}
\providecommand{\covolMah}{\operatorname{m}}
\providecommand{\covolFs}{\lambda}
\providecommand{\covolCyc}{\operatorname{I}}
\providecommand{\covolVt}{\operatorname{vt}}
\providecommand{\covolPt}{\mathfrak{p}}
\providecommand{\covolRel}{\operatorname{Rel}}
\providecommand{\covolLm}{\operatorname{e}}
\providecommand{\covolCdeg}{\operatorname{deg}_{\mathrm{Ch}}}
\providecommand{\covolHar}{\mathrm{H}}
\providecommand{\covolSmall}{\mathcal{S}}
\providecommand{\covolGen}{\operatorname{gv}}
\providecommand{\covolImDeg}{\delta_{\mathrm{im}}}
\providecommand{\covolImHt}{h_{\mathrm{im}}}

This section proves \cref{thm:covol} and \cref{thm:covolume:assembly}. The formal proof of
\cref{thm:covol} is an assembly: it applies the lift estimates (\cref{thm:covolume:lift}) to the
mean-slope statement $\sfname{DeepStR}$ (\cref{thm:deepstr}) and to four packages of hypotheses, which
it builds from the fifteen inputs of \cref{tab:covolume:inputs} and from four results proved in this
section. These four results carry the mathematical content of the section.
\begin{enumerate}[label=(\arabic*)]
\item \emph{The Chow height formula} (\cref{pf:covol:thm:chowformheight}). Let $X$ be an arithmetic
  variety over the integers $\OO_K$ of a number field, embedded in projective space by sections
  $s_1,\dots,s_m$ of a line bundle, and give the line bundle the Fubini--Study metrics attached to
  Hermitian norms on $\C^m$. Then its arithmetic self-intersection is $[K:\Q]$ times the height of a
  Chow form of $X_K$, measured by Gauss norms at the finite places and by Mahler measures on products
  of spheres at the infinite places, plus an explicit harmonic correction. This is the comparison of
  heights of varieties and heights of Chow forms of Philippon and Soul\'e
  \cite{Philippon1991,Soule1991}, see also \cite[\S4.3]{BostGilletSoule1994}; the development proves
  it from the axioms of arithmetic heights by induction on the dimension, and we give this induction in
  full (\cref{pf:covol:thm:main}).
\item \emph{The Chow form height comparison} (\cref{pf:covol:thm:comparison}, the second root of this
  section). For a Hermitian lattice $W$, a linear map $\mathrm{ev}$ from $KW$ to a field whose ratios
  generate it, and a normal model $(X,\overline L)$ of the image of $\mathrm{ev}$ on which the sections
  of $W$ carry the Fubini--Study metric of the lattice, the arithmetic self-intersection of
  $\overline L(t)$ equals $[K:\Q](D+1)\,\delta\,(h+t)$, where $D$, $\delta$ and $h$ are the dimension,
  the degree and the normalised Chow height of the image. The formal statement is the corresponding
  inequality.
\item \emph{Injectivity of small forms} (\cref{pf:covol:thm:small}). On a good deep cover, small cusp
  forms of weight $8i\cdot u$ give distinct small sections of $\overline L(t)^{\otimes 8i}$ on any such
  model.
\item \emph{Non-vanishing at interior points} (\cref{pf:covol:thm:nv}). At every complex point of a
  unitary moduli space and for every dominant weight $a$ there are a dominant weight $a'$ and
  algebraic automorphic forms of weights $a'$ and $a+a'$ which do not vanish at the point.
\end{enumerate}
The second result is the \emph{self-intersection hypothesis} of the lift estimates, and the third,
applied with the fourth, is their \emph{small-sections hypothesis} (\cref{pf:covol:def:packages}); the lift
estimates use them to compare numbers of small automorphic forms with arithmetic volumes.

\Cref{sec:pf:covol:heights} recalls the formal notions of arithmetic heights and Chow heights.
\Cref{sec:pf:covol:chowht} proves the Chow height formula, \cref{sec:pf:covol:comparison} the Chow form
height comparison, \cref{sec:pf:covol:small} the injectivity of small forms and \cref{sec:pf:covol:nv}
the non-vanishing statement. \Cref{sec:pf:covol:literature} collects the statements from the literature
used by the assembly, \cref{sec:pf:covol:assembly} proves \cref{thm:covolume:assembly} and
\cref{sec:pf:covol:main} proves \cref{thm:covol}; \cref{sec:pf:covol:further} proves further results of the
development used in other sections of Part~II. Many auxiliary statements of the formal proof are
classical facts about Chow forms, $\mathrm{Proj}$ and Mahler measures; we state them precisely and give
references to the literature together with the formal statements which prove them.

Throughout, $K$ is a number field, $\Sigma_K$ is the set of its $[K:\Q]$ embeddings $\sigma\colon K\to\C$,
and for a maximal ideal $\mathfrak q$ of $\OO_K$ we write $|\cdot|_{\mathfrak q}$ for the absolute value
with $|a|_{\mathfrak q}=\Nm(\mathfrak q)^{-\ord_{\mathfrak q}(a)}$. For a polynomial $f$ over $K$,
$|f|_{\mathfrak q}$ is the maximum of $|c|_{\mathfrak q}$ over the coefficients $c$ of $f$ (the Gauss
norm). $\covolHar_n=\sum_{i=1}^n\frac1i$ is the $n$-th harmonic number ($\covolHar_0=0$). We write
$\mathcal P$ for the Hodge package of \cref{tab:covolume:inputs}. Among its statements are three facts
of arithmetic intersection theory used below: the \emph{existence of arithmetic height functions} for
all Hermitian line bundles on arithmetic varieties (\cref{pf:covol:def:height}), the \emph{projection
formula} (if a morphism of arithmetic varieties of the same dimension maps the generic point to the
generic point with residue degree $r\ne0$, the pull-back of a Hermitian line bundle with
self-intersection $a$ has self-intersection $ra$); and it implies the formula for Fubini--Study heights of
subvarieties of projective space recalled in \cref{pf:covol:lit:fs}.

% ---------------------------------------------------------------------------
\subsection{Arithmetic heights and Chow heights}\label{sec:pf:covol:heights}

\begin{definition}[Arithmetic varieties and Hermitian line bundles]\label{pf:covol:def:hermlb}
Let $R$ be a commutative ring. A morphism $p\colon X\to\Spec R$ is an \emph{arithmetic variety of
relative dimension $D$} if $R$ is a domain which is finite over $\Z$, $X$ is integral, $p$ is
projective, $X$ is flat over $\Z$ and $\dim X=D+1$. A \emph{Hermitian line bundle} $\overline L$ on
$X$ is a line bundle $L$, given by trivialisations over an open cover $(U_i)$ with transition functions
$g_{ij}$, together with, for every ring homomorphism $\sigma\colon R\to\C$, smooth functions
$\varphi_{\sigma,i}$ (the \emph{weights}) on the complex points of $U_i$ over $\sigma$, such that
$\varphi_{\sigma,j}=\varphi_{\sigma,i}-\log|g_{ij}|^2$ and the weights at $\sigma$ and $\bar\sigma$
correspond under complex conjugation. The norm of a section $s$ with local components $s_i$ at a complex
point $y$ is $\|s(y)\|^2=|s_i(y)|^2e^{-\varphi_{\sigma,i}(y)}$. For $t\in\R$, $\overline L(t)$ is
$\overline L$ with the weights $\varphi_{\sigma,i}+2t$ (norms multiplied by $e^{-t}$), and
$\overline L^{\otimes m}$ has the weights $m\varphi_{\sigma,i}$. A global section $s$ is \emph{small} if
$\|s(y)\|\le1$ at all complex points $y$ over all $\sigma$.
\formal{GlobalTheory.IsArithVar}\formal{GlobalTheory.HermLB}\formal{GlobalTheory.HermLB.smallSec}
\end{definition}

\begin{definition}[Arithmetic heights and self-intersections]\label{pf:covol:def:height}
Let $p\colon X\to\Spec R$ be an arithmetic variety and $\overline L$ a Hermitian line bundle on $X$.
Every point of $X$ is the generic point $\eta_Z$ of a unique integral closed subscheme $Z$. A function
$h\colon X\to\R$ is an \emph{arithmetic height function} for $\overline L$ if for every integral closed
subscheme $Z\subseteq X$:
\begin{enumerate}[label=(\Alph*)]
\item[(V)] if $\Gamma(Z,\OO_Z)$ is an algebra over a finite field $\kappa$ and $\dim Z=d$, then
  $h(\eta_Z)=\deg_\kappa\bigl(c_1(L|_Z)^d\bigr)\log\#\kappa$, where $\deg_\kappa(c_1(L|_Z)^d)$ is the
  top intersection number of $L|_Z$ over $\kappa$;
\item[(P)] if $\Gamma(Z,\OO_Z)$ has characteristic $0$ and $\dim Z=1$, then for every rational section
  $r$ of $L|_Z$
  \[
  h(\eta_Z)=\sum_{z}\ord_z(r)\log\#\kappa(z)-\sum_{\sigma}\sum_{y\in Z_\sigma(\C)}\log\|r(y)\|;
  \]
\item[(R)] if $\Gamma(Z,\OO_Z)$ has characteristic $0$ and $\dim Z=e+1\ge2$, then for every $m\ge0$
  and every rational section $r$ of $L^{\otimes m}|_Z$ the function $\log\|r\|$ is integrable for the
  Chern measures $c_1(\overline L|_Z)^{e}$ on the sets $Z_\sigma(\C)\smallsetminus|\mathrm{div}(r)|$, and
  \begin{equation}\label{pf:covol:eq:recursion}
  m\,h(\eta_Z)=\sum_z\ord_z(r)\,h(z)-\sum_\sigma\int_{Z_\sigma(\C)\smallsetminus|\mathrm{div}(r)|}
  \log\|r\|\,c_1(\overline L|_Z)^{e}.
  \end{equation}
\end{enumerate}
Here $z$ runs over the points of codimension one of $Z$ and $\ord_z(r)$ is the order of the Cartier
divisor of $r$ at $z$. We say that $\overline L$ \emph{has self-intersection $a$} if there is an
arithmetic height function $h$ with $h(\eta_X)=a$.
\formallabel{GlobalTheory.HermLB.IsArithHeight}\formal{GlobalTheory.HermLB.IsSelfInter}
\end{definition}

The axioms (P) and (R) are the inductive formula of Bost, Gillet and Soul\'e for heights of cycles
\cite[\S3]{BostGilletSoule1994}, and (V) computes the height of a cycle in a fibre at a finite
prime. The existence of arithmetic height functions for all Hermitian line bundles on arithmetic
varieties is part of the Hodge package $\mathcal P$; it is the existence of the heights of
\cite{GilletSoule1990,BostGilletSoule1994}. Uniqueness is formal:

\begin{lemma}[Uniqueness of heights]\label{pf:covol:lem:unique}
Two arithmetic height functions for the same Hermitian line bundle on an arithmetic variety coincide.
In particular a Hermitian line bundle has at most one self-intersection number, and isometric Hermitian
line bundles have the same arithmetic height functions.
\formallabel{GlobalTheory.HermLB.IsArithHeight.unique}
\end{lemma}

\begin{proof}
Let $h,h'$ be arithmetic height functions. $X$ is Noetherian, so we may argue by Noetherian induction on
the closure of a point $x=\eta_Z$. Every integral closed $Z$ is either of positive characteristic $\ell$,
in which case $\Gamma(Z,\OO_Z)$ is an algebra over $\F_\ell$ and (V) gives the same value for $h$ and $h'$;
or $\Gamma(Z,\OO_Z)$ has characteristic $0$ and $\dim Z=1$, in which case (P) gives the same value;
or $\Gamma(Z,\OO_Z)$ has characteristic $0$ and $\dim Z=e+1\ge2$. In the last case choose a rational
section $r$ of $L|_Z$ and apply (R) with $m=1$: the right-hand side involves only the values at points
$z$ with $\ord_z(r)\ne0$, whose closures are strictly smaller than $Z$, and the same archimedean
integral. The last assertion holds because the axioms involve the metrics only through the norms
$\|r\|$.
\end{proof}

\begin{definition}[Embeddings by sections; point ideals; Fubini--Study metrics]\label{pf:covol:def:emb}
Let $p\colon Z\to\Spec R$, $B$ a line bundle on $Z$ and $s_1,\dots,s_m\in\Gamma(Z,B)$.
\begin{enumerate}[label=(\alph*)]
\item The sections \emph{embed} $Z$ if they have no common zero, every non-vanishing locus $D(s_k)$ is
  affine, and $\Gamma(D(s_k),\OO_Z)$ is generated as an $R$-algebra by the ratios $s_j/s_k$. (Then the
  $s_k$ define a closed immersion of $Z$ into $\mathbb P^{m-1}_R$ under which $B$ is the restriction of
  $\OO(1)$.)
\item The \emph{point ideal} $\covolPt(z)\subseteq R[X_1,\dots,X_m]$ of $z\in Z$ is the kernel of the
  homomorphism $R[X_1,\dots,X_m]\to\kappa(z)[T]$, $X_k\mapsto T\,s_k(z)$, where $s_k(z)$ is the value of
  $s_k$ at $z$ in a trivialisation of $B$ near $z$; it does not depend on the trivialisation. If
  $R=\OO_K$ and $Z$ is integral, the \emph{generic ideal} of $Z$ is $I_Z=\covolPt(\eta_Z)\otimes K
  \subseteq K[X_1,\dots,X_m]$. For $u\in R^m$ we write $s_u=\sum_ku_ks_k$ and
  $\ell_u=\sum_ku_kX_k$.
\item Let $R=\OO_K$. A \emph{Hermitian norm} on $\C^m$ is given by an invertible matrix $M$; its norm
  is $|x|_M=|Mx|$ and its dual norm is $|y|^*_M=|(M^{\mathsf T})^{-1}y|$, where $|\cdot|$ is the
  standard Euclidean norm. A weight $\varphi$ on $B$ at $\sigma$ is the \emph{Fubini--Study weight of
  $(s,M)$} if $\varphi_i(y)=\log\bigl(|(s_1(y),\dots,s_m(y))|^*_M\bigr)^2$ in every trivialisation $i$,
  so that $\|s_u(y)\|=|\ell_{\sigma u}(s(y))|/|s(y)|^*_M$.
\end{enumerate}
\formallabel{ChowFormHeightComparison₂.CF.Emb}\formallabel{ChowFormHeightComparison₂.CF.ptIdeal}%
\formallabel{ChowFormHeightComparison₂.CF.IsFS}
\end{definition}

\begin{definition}[Chow forms]\label{pf:covol:def:chow}
Let $F$ be a field and $n,k\in\N$. Polynomials in $k$ blocks of $n$ variables $u=(u_{i,j})_{i<k,j<n}$
are written $f(u_0,\dots,u_{k-1})$; $f$ is \emph{multihomogeneous of degree $\delta$} if it is
homogeneous of degree $\delta$ in each block. For $u\in F^n$, $f(u,\cdot)$ is the polynomial in $k-1$
blocks obtained by substituting $u$ for the first block. A \emph{projective prime} is a homogeneous prime
ideal $I\subseteq F[X_1,\dots,X_n]$ not containing all the $X_j$.
\begin{enumerate}[label=(\alph*)]
\item The \emph{incidence set} of an ideal $I$ is the set of $u\in\overline F^{k\times n}$ for which
  there is $x\ne0$ in $\overline F^n$ with $g(x)=0$ for all $g\in I$ and $\sum_ju_{i,j}x_j=0$ for all
  $i$. A \emph{Chow form} of $I$ in $k$ blocks is a nonzero $f$ such that a polynomial $g$ vanishes on
  the incidence set if and only if $f$ divides $g$. The \emph{Chow degree} $\covolCdeg(I,k)$ is the least
  $\delta$ for which $I$ has a multihomogeneous Chow form of degree $\delta$ in $k$ blocks.
\item An \emph{integral Chow form} of $I$ of dimension $e$ and degree $\delta$ is an irreducible
  polynomial $f$ in $e+1$ blocks, multihomogeneous of degree $\delta$, whose zero set in
  $\overline F^{(e+1)\times n}$ is the incidence set of $I$.
\item For an ideal $J$ in a Noetherian polynomial ring, $\covolRel(J)$ is the set of minimal primes of
  $J$ not containing all $X_j$, and the \emph{fundamental cycle} of $J$ is
  $[J]=\sum_{W\in\covolRel(J)}\covolLm(J,W)\,W$, where $\covolLm(J,W)$ is the length of the
  localisation at $W$ of the quotient by $J$. A \emph{Chow form of a cycle} $\sum_Wm_WW$ of projective
  primes in $k$ blocks is a polynomial $c\prod_Wg_W^{m_W}$ with $c\in F^\times$ and $g_W$ a Chow form of
  $W$ in $k$ blocks.
\end{enumerate}
\formallabel{GlobalTheory.ChowHt.IsChowFormInt}\formal{GroundTheory.Heights₂.Heights.IsChowForm}
\end{definition}

An integral Chow form of a projective prime of dimension $e+1$ (Krull dimension of the quotient) is
a Chow form in $e+1$ blocks: its zero set is the incidence hypersurface, and it is irreducible, so the
Nullstellensatz applies. These are the classical Chow forms (associated forms) of Chow and van der
Waerden \cite{ChowVanDerWaerden1937}; their existence and uniqueness up to scalars is recalled in
\cref{pf:covol:lit:chowint}.

\begin{proposition}[Degrees of integral Chow forms]\label{pf:covol:prop:hilbdeg}
Let $k$ be a field of characteristic $0$, $P\subseteq k[X_1,\dots,X_N]$ a homogeneous prime and $f$ an
integral Chow form of $P$ of dimension $D$ and degree $\delta$. Then $\delta$ is the degree of $P$, that is,
the Hilbert polynomial $H_P$ of $k[X]/P$ has degree $D$ and $\delta=D!\,a_D$, where $a_D$ is its leading
coefficient.
\formallabel{GlobalTheory.ChowExists.degree_eq_hilbertDegree}
\end{proposition}

\begin{proof}
The quotient $k[X]/P$ has dimension $D+1$ (\lref{GlobalTheory.ChowExists.ringKrullDim_of_isChowFormInt},
which derives this from the incidence description of the zero set of $f$), so $H_P$ has degree $D$ and the
degree of $P$ is $D!$ times its coefficient of degree $D$
(\lref{GlobalTheory.degree_idealHilbertPoly_eq_of_isPrime},
\lref{GlobalTheory.cast_hilbertDegree_of_degree_eq}; \cite[Thm.~I.7.5]{Hartshorne1977}). We prove
$\delta=D!\,a_D(H_P)$, where $a_D(H)$ denotes the coefficient of degree $D$, by induction on $D$, for all
fields of characteristic $0$ simultaneously.

\emph{$D=0$.} The zero set of $P$ is a single line: there are $y\in\bar k^N$ and $j$ with $y_j=1$ such that
$P$ is the ideal of polynomials vanishing at $y$ (\lref{GlobalTheory.ChowExists.exists_eq_coneIdeal}).
Then $H_P$ is the constant $[k(y):k]$, where $k(y)$ is generated by the coordinates of $y$
(\lref{GlobalTheory.CompHt.idealHilbertPoly_coneIdeal_eq_C_finrank_adjoin}), and the degree of an
integral Chow form of $P$ in one block is $[k(y):k]$: up to a scalar it is the product of the linear forms
$\ell_u(y')$ over the conjugates $y'$ of $y$ (\lref{GlobalTheory.CompHt.PointUp.le_finrank_of_isChowFormInt},
\lref{GlobalTheory.CompHt.PointLow.finrank_le_of_isChowFormInt}).

\emph{$D\to D+1$.} Let $L=k(T_1,\dots,T_N)$ be a rational function field and $\ell_T=\sum_jT_jX_j$ the generic
linear form; let $P_L=PL[X]$. Then:
(a) $H_{P_L}=H_P$, as the graded pieces of $k[X]/P$ and $L[X]/P_L$ have the same dimensions
(\lref{GlobalTheory.HilbBaseChange.idealHilbertPoly_map});
(b) $P_L$ is prime and $\ell_T\notin P_L$ (\lref{GlobalTheory.ExtPrime.isPrime_map_and_generic_form_notMem};
some $X_j$ is not in $P$, since $f$ is a Chow form);
(c) there is a prime $Q\supseteq P_L+(\ell_T)$ which is the saturation of $P_L+(\ell_T)$ with respect to the
variables, i.e. for every $g\in Q$ and every $j$ some $X_j^mg$ lies in $P_L+(\ell_T)$
(\lref{GlobalTheory.Chow.exists_prime_saturation_generic_section}; its hypothesis, that for every $u$ the
zero set of $P$ meets the hyperplane $\ell_u=0$ outside $0$, holds because $f$ vanishes whenever one block
is $0$ and its zero set is the incidence set);
(d) the generic specialisation of $f$ gives an integral Chow form $f'$ of $P_L+(\ell_T)$, hence of $Q$, of
dimension $D$ and degree $\delta$ (\lref{GlobalTheory.ChowSpec.exists_isChowFormInt_generic_section};
the passage to $Q$ uses that $Q$ and $P_L+(\ell_T)$ have the same zeros outside $0$, by the saturation
property);
(e) $Q$ is homogeneous and $H_Q=H_{P_L+(\ell_T)}$
(\lref{GlobalTheory.prime_saturation_homogeneous_and_hilbertPoly_eq});
(f) $H_{P+(\ell)}(t)=H_P(t)-H_P(t-1)$ for a homogeneous prime $P$ and a linear form $\ell\notin P$, by the
exact sequence $0\to(k[X]/P)(-1)\xrightarrow{\ell}k[X]/P\to k[X]/(P+(\ell))\to0$
(\lref{GlobalTheory.idealHilbertPoly_sup_of_isPrime}).
If $\deg H\le D+1$, the coefficient of degree $D$ of $H(t)-H(t-1)$ is $(D+1)\,a_{D+1}(H)$. By the induction
hypothesis for $L$, $Q$ and $f'$, and since $\deg H_P=D+1$,
\[
\delta=D!\,a_D(H_Q)=D!\,a_D\bigl(H_P(t)-H_P(t-1)\bigr)=D!\,(D+1)\,a_{D+1}(H_P)=(D+1)!\,a_{D+1}(H_P).\qedhere
\]
\end{proof}

\begin{definition}[Mahler measures, cone integrals and Chow heights]\label{pf:covol:def:chowht}
Let $M$ define a Hermitian norm on $\C^n$ and $k\in\N$.
\begin{enumerate}[label=(\alph*)]
\item The \emph{Mahler measure} of a polynomial $f$ in $k$ blocks of $n$ variables over $\C$ is
  \[
  \covolMah_M(f)=\int\log|f(z_0,\dots,z_{k-1})|\,d\mu(z),
  \]
  where $\mu$ is the product of the normalised uniform measures on the $M$-unit spheres
  $\{|z_i|_M=1\}$. (Formally, it is the average over the product of the unit balls of
  $\log|f(z_0/|z_0|_M,\dots)|$, which is the same.)
\item For $u\in\C^n$, $\covolFs_M(u;y)=\log\bigl(|\ell_u(y)|/|y|^*_M\bigr)$.
\item For a homogeneous ideal $W\subseteq\C[X_1,\dots,X_n]$, $e\in\N$ and a function $g$ on $\C^n$,
  the \emph{cone integral} $\covolCyc_M(W,e,g)$ is the integral of $g(M^{\mathsf T}y)$ over the set of
  $y$ with $|y|<1$ and $M^{\mathsf T}y$ in the zero set of $W$, for the $2(e+1)$-dimensional Hausdorff
  measure, divided by the volume of the unit ball of $\C^{e+1}$. For a cycle $\sum m_WW$,
  $\covolCyc_M(\sum m_WW,e,g)=\sum m_W\covolCyc_M(W,e,g)$.
\item \emph{Archimedean data} $H=(M_\sigma)_{\sigma\in\Sigma_K}$ on $\C^n$ consist of Hermitian norms
  with $|\bar u|_{M_{\bar\sigma}}=|u|_{M_\sigma}$. The \emph{Chow height} of a polynomial $f$ over $K$
  in $k$ blocks of $n$ variables, for the degree $\delta$, is
  \begin{equation}\label{pf:covol:eq:chowht}
  \covolHch^H(k,\delta,f)=\frac1{[K:\Q]}\Bigl(\sum_{\mathfrak q}\log|f|_{\mathfrak q}+
  \sum_{\sigma\in\Sigma_K}\Bigl(\covolMah_{M_\sigma}(\sigma f)+\frac{k\,\delta\,\covolHar_{n-1}}2\Bigr)\Bigr).
  \end{equation}
\item For a projective prime $I\subseteq K[X_1,\dots,X_n]$, $u\in\OO_K^n$, $k\in\N$ and a maximal ideal
  $\mathfrak q$, let $\mathcal I=I\cap\OO_K[X]$, let $\mathcal V_{\mathfrak q}$ be the set of
  $W\in\covolRel(\mathcal I+(\ell_u))$ with $W\cap\OO_K=\mathfrak q$, and
  \[
  \covolVt_{\mathfrak q}(I,k,u)=\sum_{W\in\mathcal V_{\mathfrak q}}\covolLm\bigl(\mathcal I+(\ell_u),W\bigr)\,
  \covolCdeg(W\bmod\mathfrak q,k)\log\Nm(\mathfrak q),
  \]
  where $W\bmod\mathfrak q$ is the image of $W$ in $(\OO_K/\mathfrak q)[X]$, and
  $\covolVt(I,k,u)=\sum_{\mathfrak q}\covolVt_{\mathfrak q}(I,k,u)$.
\end{enumerate}
\formallabel{GroundTheory.Heights₂.Heights.mahler}\formallabel{GroundTheory.Heights₂.Heights.cycleIntegral}%
\formallabel{GroundTheory.Heights₂.Heights.chowHt}\formallabel{GroundTheory.Heights₂.Heights.vertTerm}
\end{definition}

For the standard norm and a unit vector $p\in\C^n$ one has $\covolMah_1(\ell_p)=-\frac12\covolHar_{n-1}$:
by unitary invariance we may take $p=(1,0,\dots,0)$, and $|z_1|^2$ has the beta distribution with
parameters $(1,n-1)$ on the unit sphere, so its mean logarithm is $\psi(1)-\psi(n)=-\covolHar_{n-1}$, where
$\psi$ is the digamma function. Hence for a point $x\in\mathbb P^{n-1}(K)$, whose Chow form in one block is
$\ell_x$, the Chow height for the standard norms is the Weil height
$\frac1{[K:\Q]}\bigl(\sum_{\mathfrak q}\log|x|_{\mathfrak q}+\sum_\sigma\log|\sigma x|\bigr)$ with $L^2$ norms
at the infinite places. In general the harmonic correction makes $\covolHch$ the height of the variety for
the Fubini--Study metric; see \cref{pf:covol:rem:classical}.

The second normalisation of the same quantity used in the development is the following.

\begin{definition}[Form heights]\label{pf:covol:def:formht}
Let $w\in K^{N\times n}$, let $A=(A_\sigma)_{\sigma\in\Sigma_K}$ be invertible complex $N\times N$
matrices, and let $f$ be a polynomial over $K$ in $B$ blocks of $N$ variables. Write
$f\circ w$ for the polynomial in $B$ blocks of $n$ variables $f(wu_0',\dots,wu_{B-1}')$. The
\emph{form height} of $f$ for the degree $\delta$ is
\[
\covolHform(\delta,w,A,f)=\frac1{[K:\Q]}\Bigl(\sum_{\mathfrak q}\log|f\circ w|_{\mathfrak q}+
\sum_{\sigma}\Bigl(\covolMah^{A_\sigma}(\sigma f)+\frac{B\,\delta\,\covolHar_{N-1}}2\Bigr)\Bigr),
\]
where $\covolMah^{A_\sigma}(\sigma f)=\int\log|\sigma f(A_\sigma z_0,\dots,A_\sigma z_{B-1})|\,d\mu(z)$ and
$\mu$ is the product of the uniform probability measures on the unit spheres of $\C^N$. The frames
$A$ are \emph{conjugation compatible} if $|A_{\bar\sigma}^{-1}\bar x|=|A_\sigma^{-1}x|$ for all
$\sigma$ and $x$.
\formallabel{GlobalTheory.ChowHt.formHeight}
\end{definition}

% ---------------------------------------------------------------------------
\subsection{The Chow height formula}\label{sec:pf:covol:chowht}

We first list the auxiliary statements used in the proof. The first lemma collects facts about
embeddings by sections; they are standard properties of $\mathrm{Proj}$ and of Cartier divisors
(\cite[Ch.~II]{Hartshorne1977}, \cite[Ch.~1--2]{Fulton1998}).

\begin{lemma}[Embedded arithmetic varieties]\label{pf:covol:lem:emb}
Let $R$ be a commutative ring, $p\colon Z\to\Spec R$ with $Z$ integral, and let $s_1,\dots,s_m$ embed $Z$.
\begin{enumerate}[label=(\alph*)]
\item Every $\covolPt(z)$ is a projective prime; $z$ specialises to $z'$ if and only if
  $\covolPt(z)\subseteq\covolPt(z')$; every projective prime containing $\covolPt(\eta_Z)$ is
  $\covolPt(z)$ for some $z$; and if $u\in R^m$, then $z\in D(s_u)$ if and only if
  $\ell_u\notin\covolPt(z)$. In particular $z\mapsto\covolPt(z)$ is injective.
\item If $i\colon Y\to Z$ is a closed immersion, the restricted sections embed $Y$ and
  $\covolPt_Y(y)=\covolPt_Z(i(y))$.
\item If $R$ is Noetherian, $Z$ locally Noetherian, $u\in R^m$ and $\eta_Z\in D(s_u)$, then
  $\ord_z(s_u)\ne0$ if and only if $\covolPt(z)\in\covolRel(\covolPt(\eta_Z)+(\ell_u))$, and then
  $\ord_z(s_u)=\covolLm(\covolPt(\eta_Z)+(\ell_u),\covolPt(z))$.
\item If $Y$ is an arithmetic variety of relative dimension $D$ and $Z\subseteq Y$ is an integral closed
  subscheme whose generic point has codimension one in $Y$, then $\dim Z=D$.
\end{enumerate}
\end{lemma}

\begin{proof}
(a)--(d) are \lref{ChowFormHeightComparison₂.CF.stPts}, \lref{ChowFormHeightComparison₂.CF.stPull},
\lref{ChowFormHeightComparison₂.CF.stDiv} and \lref{ChowFormHeightComparison₂.CF.stDim}. By (a)
and (b), the embedding identifies $Z$ with a closed subscheme of $\mathbb P^{m-1}_R$ and $\covolPt(z)$ with
the homogeneous ideal of the closure of $z$, which gives (a) by the correspondence between points of
$\mathrm{Proj}$ and relevant homogeneous primes; (c) is the computation of the order of a Cartier divisor
on a Noetherian integral scheme as a length, transported to the homogeneous coordinate ring; (d) is the
dimension formula for integral schemes of finite type and flat over a Dedekind domain.
\end{proof}

The second lemma collects facts about Chow forms. Parts (a), (b), (c) and (e) are classical properties of
Chow forms (\cite{ChowVanDerWaerden1937}, \cite[\S1]{Philippon1986}); (d) is the Nullstellensatz; (f) is
Gauss's lemma together with the additivity of $\log|fg|=\log|f|+\log|g|$.

\begin{lemma}[Chow forms and Mahler measures]\label{pf:covol:lem:chowcalc}
Let $F$ be a field and $I\subseteq F[X_1,\dots,X_n]$ a projective prime.
\begin{enumerate}[label=(\alph*)]
\item If $F$ has characteristic $0$, $f$ is a Chow form of $I$ in $e+1$ blocks and $u\in F^n$ with
  $\ell_u\notin I$, then $f(u,\cdot)$ is a Chow form of the cycle $[I+(\ell_u)]$ in $e$ blocks.
\item If $F=K$, $f$ is a Chow form of $I$ in $k$ blocks and $\sigma\in\Sigma_K$, then every
  component of $[\sigma I]$ is a projective prime and $\sigma f$ is a Chow form of $[\sigma I]$.
\item If $F=K$, $\ell_u\notin I$ and $\sigma\in\Sigma_K$, then $\ell_{\sigma u}$ lies in no component of
  $[\sigma I]$.
\item $I$ has a zero $x\ne0$ in $\overline F^n$. Consequently $I$ has no Chow form in $0$ blocks.
\item A divisor of a multihomogeneous polynomial is multihomogeneous.
\item For nonzero $f_1,f_2$ over $K$ in $k$ blocks and any archimedean data $H$,
  $\covolHch^H(k,\delta_1+\delta_2,f_1f_2)=\covolHch^H(k,\delta_1,f_1)+\covolHch^H(k,\delta_2,f_2)$;
  and for a nonzero constant $c$, $\covolHch^H(k,\delta,c)=\frac{k\delta}2\covolHar_{n-1}$, which vanishes
  if $k=0$ or $\delta=0$.
\end{enumerate}
\end{lemma}

\begin{proof}
(a) is \lref{GroundTheory.Heights₂.Heights.ChowFormSection.full}, (b) is
\lref{GroundTheory.Heights₂.Heights.ChowFormBaseChange.full} (the ideal generated by $\sigma(I)$ is
radical because $K$ is perfect, $\sigma f$ is a Chow form of it, and a Chow form of a radical homogeneous
ideal is a Chow form of its fundamental cycle, all multiplicities being $1$), (c) is
\lref{GroundTheory.Heights₂.Heights.BaseChangeAvoid.full}, (d) is
\lref{GroundTheory.Heights₂.Heights.ProjPrimeHasPoint.full} and the first assertion of (f) is
\lref{GroundTheory.Heights₂.Heights.ChowHeightMul.full}. In (c), the minimal primes of
$\sigma(I)\C[X]$ contract to minimal primes of $I$, that is to $I$, since $\C[X]$ is flat over $K[X]$;
so $\ell_{\sigma u}\in W$ would give $\ell_u\in I$. In (d), the zero exists by the Nullstellensatz because
$I$ is homogeneous and does not contain all $X_j$; with $0$ blocks the incidence set is then the whole
(one-point) space, so a Chow form $f$ would be a nonzero constant which divides $1$ although $1$ does not
vanish on the incidence set. (e) is elementary. In (f), the Gauss norms are multiplicative (Gauss's lemma)
and $\covolMah_M(f_1f_2)=\covolMah_M(f_1)+\covolMah_M(f_2)$ because $\log|f|$ is integrable on products of
spheres for $f\ne0$; the harmonic terms are additive in $\delta$; for a constant the non-harmonic part is
$\frac1{[K:\Q]}(\sum_{\mathfrak q}\log|c|_{\mathfrak q}+\sum_\sigma\log|\sigma c|)=0$ by the product
formula.
\end{proof}

The four remaining ingredients are proved in full in the rest of this subsection, before the induction:
\cref{pf:covol:prop:vert} computes the self-intersection of a vertical subvariety,
\cref{pf:covol:prop:arch} identifies the archimedean integrals with cone integrals (Wirtinger's theorem,
\cite{GriffithsHarris1978}),
and \cref{pf:covol:prop:csv,pf:covol:prop:mahler} compare the Gauss norms and the Mahler measures of a Chow
form with those of its specialisation at a linear form. The last two are the non-archimedean and the
archimedean halves of the comparison of Philippon and Soul\'e between heights of varieties and heights of
Chow forms \cite{Philippon1991,Soule1991}, \cite[\S4.3]{BostGilletSoule1994}.

\subsubsection*{Vertical subvarieties}
For a scheme $X$ proper over a field $\kappa$, a coherent sheaf $\mathcal F$ on $X$ and line bundles
$L_1,\dots,L_d$, the intersection number is
$(L_1\cdots L_d\cdot\mathcal F)=\sum_{S\subseteq\{1,\dots,d\}}(-1)^{|S|}\chi_\kappa\bigl(\mathcal F\otimes
\bigotimes_{j\in S}L_j^{-1}\bigr)$ (Snapper's definition, \cite[Ch.~18]{Fulton1998}); the top
self-intersection $\deg_\kappa(c_1(L)^d)$ of \cref{pf:covol:def:height} is $(L\cdots L\cdot\OO_X)$ with $d$
factors. We use two facts about Chow degrees over a perfect field $\kappa$, both classical
(\cite{ChowVanDerWaerden1937}, \cite[\S1]{Philippon1986}):
\begin{enumerate}[label=(\roman*)]
\item if $I\subseteq\kappa[X_1,\dots,X_m]$ is a projective prime with $\dim\kappa[X]/I=1$ and
  $X_k\notin I$, then $\covolCdeg(I,1)=\dim_\kappa\kappa[X]/(I+(X_k-1))$
  (\lref{GlobalTheory.CompanionTheory.Ar.Short.ChowClosing.Auto.ObjectsChow_ChowDegreeClosedPoint});
\item if $I$ is a projective prime with $\dim\kappa[X]/I=e+2$ and $\ell_u\notin I$, then every
  $W\in\covolRel(I+(\ell_u))$ is a projective prime with $\dim\kappa[X]/W=e+1$
  (\lref{GroundTheory.Heights₂.Heights.ConeDimComponent.proof_closed}), and
  $\covolCdeg(I,e+2)=\sum_{W\in\covolRel(I+(\ell_u))}\covolLm(I+(\ell_u),W)\,\covolCdeg(W,e+1)$
  (\lref{ChowFormHeightComparison₂.CF.V.chowDegreeSection}): the degree of a variety is the degree
  of a hyperplane section counted with multiplicities.
\end{enumerate}

\begin{lemma}[Degrees over a field]\label{pf:covol:lem:fieldcase}
Let $\kappa$ be a perfect field, $Z$ an integral scheme, projective over $\kappa$, embedded over $\kappa$ by
sections $s_1,\dots,s_m$ of a line bundle $B$, and $J=\covolPt(\eta_Z)\subseteq\kappa[X_1,\dots,X_m]$. If
$\dim\kappa[X]/J=e+1$, then $\dim Z=e$ and $\deg_\kappa(c_1(B)^e)=\covolCdeg(J,e+1)$.
\formallabel{ChowFormHeightComparison₂.CF.V.field_case}
\end{lemma}

\begin{proof}
First, $\dim Z+1\le\dim\kappa[X]/J$ for every $Z$ as in the lemma: by \cref{pf:covol:lem:emb}(a), a
strictly decreasing chain of irreducible closed subsets of $Z$, with generic points $z_0,\dots,z_r$, gives a
strictly increasing chain of projective primes $\covolPt(z_0)\subsetneq\dots\subsetneq\covolPt(z_r)$
containing $J$, and it can be prolonged by the prime $(X_1,\dots,X_m)$, which contains all $X_j$. We
argue by induction on $e$. Since the sections have no common zero, there is $k$ with $\eta_Z\in D(s_k)$,
so $X_k\notin J$ by \cref{pf:covol:lem:emb}(a).

\emph{$e=0$.} Then $\dim Z\le0$, so $Z$ is a point and $Z=D(s_k)$ is affine with ring generated over
$\kappa$ by the ratios $s_j/s_k$; this ring is $\kappa[X]/(J+(X_k-1))$, since the kernel of
$X_j\mapsto s_j/s_k$ is $J+(X_k-1)$ (\lref{ChowFormHeightComparison₂.CF.V.finrank_eq_of_dim_zero}). On
a zero-dimensional scheme the higher cohomology of $\OO_Z$ vanishes, so
$\deg_\kappa(c_1(B)^0)=\chi_\kappa(\OO_Z)=\dim_\kappa\Gamma(Z,\OO_Z)=\covolCdeg(J,1)$ by (i).

\emph{$e\to e+1$.} Let $u$ be the $k$-th unit vector, so that $s_u=s_k$ and $\ell_u=X_k$, and let $T$ be
the finite set of points $z$ with $\ord_z(s_k)\ne0$. By \cref{pf:covol:lem:emb}(a),(c) and (ii),
$z\mapsto\covolPt(z)$ is a bijection from $T$ onto $\covolRel(J+(X_k))$ with
$\ord_z(s_k)=\covolLm(J+(X_k),\covolPt(z))$: every $W\in\covolRel(J+(X_k))$ is a projective prime
containing $J$, hence equal to $\covolPt(z)$ for some $z$, and $\ord_z(s_k)\ne0$ by (c). For $z\in T$ the
closure $Y_z$ of $z$ is embedded by the restricted sections with the same point ideals
(\cref{pf:covol:lem:emb}(b)) and $\dim\kappa[X]/\covolPt(z)=e+1$ by (ii); by the induction hypothesis,
$\dim Y_z=e$ and $\deg_\kappa(c_1(B|_{Y_z})^e)=\covolCdeg(\covolPt(z),e+1)$.

Since $Z$ is integral and $s_k\ne0$ at $\eta_Z$, the section $s_k\colon B^{-1}\to\OO_Z$ is injective; let
$\OO_D$ be its cokernel. Additivity of the Euler characteristic in the exact sequences
$0\to B^{-1}\otimes N\to N\to\OO_D\otimes N\to0$ gives
$(B^{e+1}\cdot\OO_Z)=(B^e\cdot\OO_Z)-(B^e\cdot B^{-1})=(B^e\cdot\OO_D)$. The support of $\OO_D$ is the zero
set of $s_k$, and every point $z$ of it outside $T$ has a closure of dimension $<e$: by
\cref{pf:covol:lem:emb}(a), $J+(X_k)\subseteq\covolPt(z)$, so $\covolPt(z)$ contains some
$W\in\covolRel(J+(X_k))$, $W=\covolPt(z')$ with $z'\in T$, $z'\ne z$, and the closure of $z$ is a proper
closed subset of $Y_{z'}$, of dimension $e$. For a coherent sheaf $\mathcal F$ whose support consists of
the points of $T$ (with closures of dimension $e$) and of points with closures of dimension $<e$,
\[
(B^e\cdot\mathcal F)=\sum_{z\in T}\len_{\OO_{Z,z}}(\mathcal F_z)\,\deg_\kappa\bigl(c_1(B|_{Y_z})^e\bigr)
\]
(\lref{ChowFormHeightComparison₂.CF.V.interNum_eq_sum}, which proves this by d\'evissage of
$\mathcal F$, the intersection number of $e$ line bundles with a sheaf supported in dimension $<e$
being $0$). For $\mathcal F=\OO_D$, $\len(\OO_{D,z})=\ord_z(s_k)$ is the order of the Cartier divisor
of $s_k$ at $z$ (\lref{ChowFormHeightComparison₂.CF.V.length_cut}). Hence, with (ii),
\begin{multline*}
\deg_\kappa(c_1(B)^{e+1})=\sum_{z\in T}\ord_z(s_k)\,\covolCdeg(\covolPt(z),e+1)\\
=\sum_{W\in\covolRel(J+(X_k))}\covolLm(J+(X_k),W)\,\covolCdeg(W,e+1)=\covolCdeg(J,e+2).
\end{multline*}
Finally $\dim Z\le e+1$ by the first paragraph. Conversely, $J+(X_k)$ is contained in the prime
$(X_1,\dots,X_m)$, hence in a minimal prime $W$ of $J+(X_k)$ contained in it; as $X_k\notin J$ and
$\kappa[X]/J$ is a domain of dimension $e+2$, $\dim\kappa[X]/W=e+1\ge1$, so $W$ does not contain all
$X_j$ and $W=\covolPt(z)$ with $z\in T$. The closure of $z$ has dimension $e$ and is a proper subset of
$Z$ (as $X_k\in W\setminus J$), so $\dim Z\ge e+1$.
\end{proof}

\begin{proposition}[Vertical subvarieties]\label{pf:covol:prop:vert}
Let $Z$ be an integral scheme, projective over $\OO_K$, embedded by sections $s_1,\dots,s_m$ of a line
bundle $B$, with $\covolPt(\eta_Z)\cap\OO_K=\mathfrak q$ a maximal ideal and $\dim Z=d$. Then
$\Gamma(Z,\OO_Z)$ is an algebra over $\kappa=\OO_K/\mathfrak q$, and
$\deg_\kappa(c_1(B)^d)=\covolCdeg(\covolPt(\eta_Z)\bmod\mathfrak q,d+1)$.
\formallabel{ChowFormHeightComparison₂.CF.stVert}
\end{proposition}

\begin{proof}
An element $r\in\OO_K$ vanishes in $\Gamma(Z,\OO_Z)$ if and only if the constant $r$ lies in
$\covolPt(\eta_Z)$, that is $r\in\mathfrak q$; so the structure map factors through $\kappa$. Over
$\kappa$ the same sections embed $Z$, $Z$ is projective over $\kappa$ (as $\OO_K\to\kappa$ is surjective),
and the point ideal of $\eta_Z$ over $\kappa$ is the image $J$ of $\covolPt(\eta_Z)$ in $\kappa[X]$ (the
reduction map $\OO_K[X]\to\kappa[X]$ is surjective). The dimension of $\kappa[X]/J$ is at most $m$ and, by
the first paragraph of the proof of \cref{pf:covol:lem:fieldcase}, at least $d+1$; write it $n$. By
\cref{pf:covol:lem:fieldcase} ($\kappa$ is finite, hence perfect), $\dim Z=n-1$ and
$\deg_\kappa(c_1(B)^{n-1})=\covolCdeg(J,n)$. Since $\dim Z=d$, $n=d+1$.
\end{proof}

\subsubsection*{Archimedean integrals}

\begin{proposition}[Wirtinger's theorem for embedded varieties]\label{pf:covol:prop:arch}
Let $Z$ be an arithmetic variety of relative dimension $e$ over $\OO_K$, embedded by sections
$s_1,\dots,s_m$ of a line bundle $B$, whose generic ideal $I_Z$ is a projective prime with a Chow form in
$e+1$ blocks. Let $\sigma\in\Sigma_K$, let $\varphi$ be a smooth Fubini--Study weight of $(s,M)$ at
$\sigma$, and let $u\in K^m$ with $\ell_u\notin I_Z$. Then
\[
\int_{Z_\sigma(\C)\cap\{\ell_{\sigma u}(s)\ne0\}}\covolFs_M(\sigma u;s(y))\,c_1(\varphi)^e=
\covolCyc_M\bigl([\sigma I_Z],e,\covolFs_M(\sigma u;\cdot)\bigr),
\]
where $c_1(\varphi)^e$ is the Chern measure and $[\sigma I_Z]$ is the fundamental cycle of the ideal of
$\C[X]$ generated by $\sigma(I_Z)$; and if $e=0$, then also
$\sum_{y\in Z_\sigma(\C)}\covolFs_M(\sigma u;s(y))=\covolCyc_M([\sigma I_Z],0,\covolFs_M(\sigma u;\cdot))$.
\formallabel{ChowFormHeightComparison₂.CF.stArch}
\end{proposition}

Here the Chern measure is defined through local holomorphic charts given by algebraic functions: on such
a chart it has the density $\det(\partial\bar\partial\varphi)$ (suitably normalised), and the integral of
$g$ over a set $S$ is $\int_Sg\,d\mu_+-\int_Sg\,d\mu_-$, where $\mu_\pm$ are the measures with the positive and
negative parts of this density.

\begin{proof}
Let $J\subseteq\C[X]$ be the ideal generated by $\sigma(I_Z)$, let $\mathcal C\subseteq\C^m$ be the set of
$y$ with $|y|<1$ and $M^{\mathsf T}y$ in the zero set of $J$, let $\mathcal H$ be the
$2(e+1)$-dimensional Hausdorff measure and $b=\mathcal H(\{y\in\C^{e+1}:|y|<1\})\in(0,\infty)$. We use four
facts.
\begin{enumerate}[label=(\arabic*)]
\item There is a homogeneous $q\in\OO_K[X]$ with $q\notin\covolPt(\eta_Z)$ such that every complex point
  $y$ of $Z_\sigma$ with $q(s(y))\ne0$ lies in the domain of an algebraic chart of dimension $e$
  (\lref{ChowFormHeightComparison₂.CF.stCharts}).
\item $e+1\le m$; the function $y\mapsto\covolFs_M(\sigma u;M^{\mathsf T}y)$ is $\mathcal H$-integrable on
  $\mathcal C$; $\covolCyc_M([J],e,g)=\covolCyc_M(J,e,g)$ for this function (all components of $J$ have
  multiplicity $1$); and for $q\notin I_Z$ the zero set of $\sigma q(M^{\mathsf T}\cdot)$ in $\mathcal C$ is
  $\mathcal H$-null (\lref{ChowFormHeightComparison₂.CF.C.cone_side}).
\item (Wirtinger) If $q$ is as in (1) and $\{0\}\cup(\mathcal C\cap\{\sigma q(M^{\mathsf T}\cdot)=0\})$ is
  $\mathcal H$-null, then for every measurable $F\colon\C^m\to[0,\infty]$ with $F(ty)=F(y)$ for
  $t\in\C^\times$,
  \[
  \int_{\mathcal C}F\,d\mathcal H=b\int_{Z_\sigma(\C)}F\bigl(\bar s(y)\bigr)\,d\mu_+(y),
  \]
  where $\bar s(y)\in\C^m$ is any representative of $s(y)$ in the coordinates of $M$
  (\lref{ChowFormHeightComparison₂.CF.C.lintegral_cone}, proved chart by chart: on a chart the cone
  is fibred over the chart by complex lines, and the coarea formula identifies the fibre integrals with
  the Fubini--Study density).
\item For a Fubini--Study weight the density is nonnegative, so $\mu_-=0$
  (\lref{ChowFormHeightComparison₂.CF.C.chernPart_neg}).
\end{enumerate}
The function $F_G=G\circ\covolFs_M(\sigma u;M^{\mathsf T}\cdot)$ is measurable and invariant under $\C^\times$
for every measurable $G\colon\R\to[0,\infty]$, because $|\ell_{\sigma u}(ty)|=|t||\ell_{\sigma u}(y)|$ and
$|ty|^*_M=|t||y|^*_M$. By (2), the hypothesis of (3) holds for the $q$ of (1), and (3) gives
$\int_{\mathcal C}G(f)\,d\mathcal H=b\int G(g)\,d\mu_+$ for all such $G$, with $f=\covolFs_M(\sigma
u;M^{\mathsf T}\cdot)$ and $g(y)=\covolFs_M(\sigma u;s(y))$. Taking $G(x)=|x|$ shows that $g$ is
$\mu_+$-integrable, as $f$ is $\mathcal H$-integrable by (2); taking for $G$ the positive and the negative
parts gives $\int g\,d\mu_+=b^{-1}\int_{\mathcal C}f\,d\mathcal H=\covolCyc_M(J,e,\covolFs_M(\sigma
u;\cdot))=\covolCyc_M([J],e,\covolFs_M(\sigma u;\cdot))$ by (2). By (4) the left-hand side is the integral
against $c_1(\varphi)^e$; and the integral over the set where $\ell_{\sigma u}(s(y))\ne0$ is the integral
over $Z_\sigma(\C)$, because on the complement the integrand is replaced by $0$ (formally $\log0=0$).
If $e=0$, every complex point of $Z_\sigma$ lies over the generic point of $Z$, so $q(s(y))\ne0$ and $y$
lies in a chart of dimension $0$ by (1); thus $Z_\sigma(\C)$ is finite and $c_1(\varphi)^0$ is the counting
measure (\lref{ChowFormHeightComparison₂.CF.C.integral_zero}), which gives the second identity.
\end{proof}

\subsubsection*{The vertical term of a Chow form}
Let $I\subseteq K[X_1,\dots,X_n]$ be a projective prime with a Chow form $f$ in $e+1$ blocks, let
$u\in\OO_K^n$ with $\ell_u\notin I$, and let $\mathfrak q$ be a maximal ideal of $\OO_K$. Write
$\mathcal I=I\cap\OO_K[X]$ and let $S_{\mathfrak q}=\covolRel(\mathcal I+\mathfrak q\OO_K[X])$ be the set of
components of the fibre of $\mathcal I$ at $\mathfrak q$. Let $t=(t_{ij})_{i\le e,j\le n}$ and
$t'=(t'_j)_{j\le n}$ be indeterminates, $R$ the localisation of $\OO_K[t,t']$ at the prime
$\mathfrak q\OO_K[t,t']$, and
\[
\mathcal B=R[X_1,\dots,X_n]\big/\bigl(\mathcal IR[X]+(\textstyle\sum_jt_{ij}X_j)_{i\le e}+
(\sum_jt'_jX_j-1)\bigr),
\]
the coordinate ring of the intersection of the model of $I$ with $e$ generic hyperplanes in a generic
affine chart. For $W\in S_{\mathfrak q}$ let $\mathfrak n_W$ be the ideal of $\mathcal B$ generated by $W$.
The proof of the next proposition uses the following statements of the development.
\begin{enumerate}[label=(\arabic*)]
\item $R$ is a discrete valuation ring, $\mathcal B$ is a finite free $R$-module, and multiplication by
  $\ell_u$ on $\mathcal B$ is injective
  (\lref{GroundTheory.Heights₂.Heights.Vertical.ParDVRIsDVR.full},
  \lref{GroundTheory.Heights₂.Heights.Vertical.SecBFinite.full},
  \lref{GroundTheory.Heights₂.Heights.Vertical.SecBTorsionFree.full},
  \lref{GroundTheory.Heights₂.Heights.Vertical.SecBLinInj.proof}).
\item \emph{Norm identity.} If $g\in\OO_K[u_0,\dots,u_e]$ is a Chow form of $I$ with a coefficient not in
  $\mathfrak q$, then $g(u,t)=w\cdot\Nm_{\mathcal B/R}(\ell_u)$ for a unit $w$ of $R$
  (\lref{GroundTheory.Heights₂.Heights.Vertical.NormIdentity.proof}); this rests on the classical
  description of a Chow form as the norm form of the generic linear form on the generic zero-dimensional
  section.
\item $W\mapsto\mathfrak n_W$ is a bijection from $S_{\mathfrak q}$ onto the set of maximal ideals of
  $\mathcal B$ (\lref{GroundTheory.Heights₂.Heights.Vertical.SpecMax.proof}), and
  $\len_R(\mathcal B/\mathfrak n_W)=\covolCdeg(W\bmod\mathfrak q,e+1)$
  (\lref{GroundTheory.Heights₂.Heights.Vertical.SpecResidue.proof}).
\item The length of $\mathcal B_{\mathfrak n_W}/\ell_u\mathcal B_{\mathfrak n_W}$ is
  $\covolLm(\mathcal I+(\ell_u),W)$ if $\ell_u\in W$ and $0$ otherwise
  (\lref{GroundTheory.Heights₂.Heights.Vertical.SpecLength.proof}); and the set
  $\mathcal V_{\mathfrak q}$ of \cref{pf:covol:def:chowht}(e) is $\{W\in S_{\mathfrak q}:\ell_u\in W\}$
  (\lref{GroundTheory.Heights₂.Heights.Vertical.VertSpec.full}).
\item For $0\ne h\in\OO_K[t,t']$, $|h|_{\mathfrak q}\Nm(\mathfrak q)^{\ord_R(h)}=1$: the valuation of $R$
  restricts to the Gauss valuation (\lref{GroundTheory.Heights₂.Heights.Vertical.GaussDVR.full}).
\end{enumerate}

\begin{proposition}[Vertical terms]\label{pf:covol:prop:csv}
In this situation,
$\log|f|_{\mathfrak q}-\log|f(u,\cdot)|_{\mathfrak q}=\covolVt_{\mathfrak q}(I,e+1,u)$.
\formallabel{GroundTheory.Heights₂.Heights.ChowFormSectionVertical.full}
\end{proposition}

\begin{proof}
Multiplying $f$ by a suitable $c\in K^\times$ we obtain $g=cf\in\OO_K[u_0,\dots,u_e]$ with a coefficient not
in $\mathfrak q$ (\lref{GroundTheory.Heights₂.Heights.Vertical.PrimitiveScale.full}); $g$ is again a
Chow form of $I$, $|g|_{\mathfrak q}=1$, and both $\log|f|_{\mathfrak q}$ and
$\log|f(u,\cdot)|_{\mathfrak q}$ change by $\log|c|_{\mathfrak q}$, so we may take $f=g$. Let
$\nu=\sum_{W\in\mathcal V_{\mathfrak q}}\covolLm(\mathcal I+(\ell_u),W)\,\covolCdeg(W\bmod\mathfrak q,e+1)$,
so that $\covolVt_{\mathfrak q}(I,e+1,u)=\nu\log\Nm(\mathfrak q)$ by definition.

For a finite free algebra $\mathcal B$ over a discrete valuation ring $R$ and a nonzerodivisor
$b\in\mathcal B$ one has $\ord_R\Nm_{\mathcal B/R}(b)=\len_R(\mathcal B/b\mathcal B)$, and
$\len_R(\mathcal B/b\mathcal B)=\sum_{\mathfrak n}\len(\mathcal B_{\mathfrak n}/b\mathcal B_{\mathfrak n})\,
\len_R(\mathcal B/\mathfrak n)$, the sum over the maximal ideals of $\mathcal B$ (standard,
\cite[App.~A]{Fulton1998}; \lref{GroundTheory.Heights₂.Heights.Vertical.NormLength.full},
\lref{GroundTheory.Heights₂.Heights.Vertical.LengthDecomp.full}). With (1)--(4),
\[
\ord_R g(u,t)=\ord_R\Nm_{\mathcal B/R}(\ell_u)=\sum_{W\in S_{\mathfrak q}}
\len\bigl(\mathcal B_{\mathfrak n_W}/\ell_u\mathcal B_{\mathfrak n_W}\bigr)\,\covolCdeg(W\bmod\mathfrak q,e+1)=\nu .
\]
The polynomial $g(u,\cdot)$ is nonzero (a Chow form of $[I+(\ell_u)]$ by \cref{pf:covol:lem:chowcalc}(a)),
so (5) gives $|g(u,\cdot)|_{\mathfrak q}=\Nm(\mathfrak q)^{-\nu}$. Hence
$\log|g|_{\mathfrak q}-\log|g(u,\cdot)|_{\mathfrak q}=\nu\log\Nm(\mathfrak q)$.
\end{proof}

\subsubsection*{The Mahler measure of a specialised Chow form}
Let $M$ define a Hermitian norm on $\C^n$. For a cycle $Z$ of projective primes of $\C[X_1,\dots,X_n]$ with
a Chow form $f$ in $e+1$ blocks and $u\in\C^n$ with $\ell_u$ in no component of $Z$, call
$D_f(u)=\covolMah_M(f(u,\cdot))-\covolCyc_M(Z,e,\covolFs_M(u;\cdot))$ the \emph{defect} of $f$ at $u$.

\begin{proposition}[Mahler measures of specialisations]\label{pf:covol:prop:mahler}
Let $Z=\sum_Wm_WW$ be a cycle of projective primes of $\C[X_1,\dots,X_n]$ with a Chow form $f$ in $e+1$
blocks which is multihomogeneous of degree $\delta$, and let $u\in\C^n$ with $\ell_u\notin W$ for all $W$ in
the support of $Z$. Then
\[
\covolMah_M(f(u,\cdot))=\covolMah_M(f)+\frac{\delta\,\covolHar_{n-1}}2+
\covolCyc_M\bigl(Z,e,\covolFs_M(u;\cdot)\bigr).
\]
\formallabel{GlobalTheory.T.mahler_section_statement_holds}
\end{proposition}

\begin{proof}
\emph{Step 1: the defect of a prime does not depend on $u$}
(\lref{GlobalTheory.T.mahler_minus_cone_integral_indep}, the formal proof of which is the following
Crofton argument). Let $W$ be a projective prime with a Chow form $g$ in $e+1$ blocks, multihomogeneous of
degree $\delta$, and let $u,u'$ with $\ell_u,\ell_{u'}\notin W$. Let $\gamma$ be the standard Gaussian
measure on $(\C^n)^e$, and for $\pi=(\pi_1,\dots,\pi_e)$ let $\hat\pi$ be the $e$ linear forms obtained by
normalising the rows of $\pi$, transported by $M$, to unit length. For $\gamma$-almost every $\pi$:
(i) the linear subspace $\{\ell_{\hat\pi_1}=\dots=\ell_{\hat\pi_e}=0\}$ meets the cone of $W$ in exactly
$\delta$ lines (\lref{GlobalTheory.T.ae_fibre_points_card_eq_degree}); (ii) $g(u,\hat\pi)\ne0$ and
$g(u',\hat\pi)\ne0$ (\lref{GlobalTheory.T.ae_polynomial_nonvanishing_rows}); and then (iii) there is
$c\ne0$ such that $g(v,\hat\pi)=c\prod_y\ell_v(y)$ for all $v$, where $y$ runs over representatives of these
lines of $M$-norm $1$ (\lref{GlobalTheory.T.chow_form_fibre_specialization}). Since
$\covolFs_M(v;y)=\log|\ell_v(y)|-\log|y|^*_M$, it follows that
$\log|g(v,\hat\pi)|-\sum_y\covolFs_M(v;y)=\log|c|+\sum_y\log|y|^*_M$ is the same for $v=u$ and $v=u'$.
Integrate over $\gamma$. The push-forward of $\gamma$ under $\pi\mapsto\hat\pi$ is a multiple $\Gamma$ of the
product of the uniform measures on the unit spheres, so
$\int\log|g(v,\hat\pi)|\,d\gamma=\Gamma\,\covolMah_M(g(v,\cdot))$
(\lref{GlobalTheory.T.sphere_mean_normalized_groups}); and by the Crofton formula
$\int\sum_y\covolFs_M(v;y)\,d\gamma(\pi)=\Gamma\,\covolCyc_M(W,e,\covolFs_M(v;\cdot))$
(\lref{GlobalTheory.T.integral_fubini_study_sum_cone}). As $0<\Gamma<\infty$
(\lref{GlobalTheory.T.gaussian_matrix_measure_pos_finite}), $D_g(u)=D_g(u')$.

\emph{Step 2: the defect of a cycle.} Write $f=c\prod_Wg_W^{m_W}$ with Chow forms $g_W$ of $W$. Since the
Mahler measure is additive on products and the cone integral is additive in the cycle,
$D_f(u)=\log|c|+\sum_Wm_WD_{g_W}(u)$ (\lref{GlobalTheory.CompanionTheory.Ar.Short.mahlerDefect_eq_sum});
by Step 1, $D_f(u)$ is the same for all admissible $u$.

\emph{Step 3: averaging.} For almost every $v$ in the $M$-unit ball, $\ell_v$ lies in no component of $Z$
(\lref{GroundTheory.Heights₂.Heights.AutoClose.AdmissibleAENrm}), so $D_f(v/|v|_M)=D_f(u)$ for almost
every $v$; the ball has positive finite volume
(\lref{GroundTheory.Heights₂.Heights.BallVolume.proof}), so $D_f(u)$ is the average of $D_f(v/|v|_M)$
over the ball. This average is $\covolMah_M(f)+\frac{\delta\covolHar_{n-1}}2$
(\lref{GroundTheory.Heights₂.Heights.AutoClose.DefectAverage}): by Fubini, the average of
$\covolMah_M(f(v,\cdot))$ is $\covolMah_M(f)$; and since the average over the unit sphere of
$\covolFs_M(v;y)$ is $-\frac12\covolHar_{n-1}$ for every $y\ne0$ (the computation after
\cref{pf:covol:def:chowht}) and $\covolCyc_M(Z,e,1)=\delta$, the average of the cone integral is
$-\frac{\delta\covolHar_{n-1}}2$. Hence $D_f(u)=\covolMah_M(f)+\frac{\delta\covolHar_{n-1}}2$, which is the
claim.
\end{proof}

\begin{proposition}[Recursion for Chow heights]\label{pf:covol:prop:recursion}
Let $I\subseteq K[X_1,\dots,X_n]$ be a projective prime with a Chow form $f$ in $e+1$ blocks which is
multihomogeneous of degree $\delta$, let $H=(M_\sigma)_\sigma$ be archimedean data and $u\in\OO_K^n$ with
$\ell_u\notin I$. Then
\begin{equation}\label{pf:covol:eq:chowrec}
\covolHch^H(e+1,\delta,f)=\covolHch^H\bigl(e,\delta,f(u,\cdot)\bigr)+\frac1{[K:\Q]}\Bigl(\covolVt(I,e+1,u)-
\sum_{\sigma\in\Sigma_K}\covolCyc_{M_\sigma}\bigl([\sigma I],e,\covolFs_{M_\sigma}(\sigma u;\cdot)\bigr)\Bigr).
\end{equation}
\formallabel{GroundTheory.Heights₂.Heights.ChowHeightRecursion.of_leaves}
\end{proposition}

\begin{proof}
Put $g=f(u,\cdot)$. Only finitely many $\log|f|_{\mathfrak q}$ and $\log|g|_{\mathfrak q}$ are nonzero,
and \cref{pf:covol:prop:csv} gives $\sum_{\mathfrak q}\log|f|_{\mathfrak q}=
\sum_{\mathfrak q}\log|g|_{\mathfrak q}+\covolVt(I,e+1,u)$. For $\sigma\in\Sigma_K$, the cycle
$[\sigma I]$ consists of projective primes, $\sigma f$ is a Chow form of it and is multihomogeneous of
degree $\delta$, and $\ell_{\sigma u}$ avoids its components (\cref{pf:covol:lem:chowcalc}(b),(c));
since $\sigma g=(\sigma f)(\sigma u,\cdot)$, \cref{pf:covol:prop:mahler} gives
$\covolMah_{M_\sigma}(\sigma f)=\covolMah_{M_\sigma}(\sigma g)-\frac{\delta\covolHar_{n-1}}2-
\covolCyc_{M_\sigma}([\sigma I],e,\covolFs_{M_\sigma}(\sigma u;\cdot))$. Inserting both identities in
\eqref{pf:covol:eq:chowht}, the harmonic terms $\frac{(e+1)\delta\covolHar_{n-1}}2$ of the left-hand side
become $\frac{e\,\delta\covolHar_{n-1}}2$ plus the term $\frac{\delta\covolHar_{n-1}}2$ just produced,
and \eqref{pf:covol:eq:chowrec} follows.
\end{proof}

We can now prove the main identity.

\begin{theorem}[Heights of embedded arithmetic varieties]\label{pf:covol:thm:main}
Let $H=(M_\sigma)_\sigma$ be archimedean data on $\C^m$, let $p\colon X\to\Spec\OO_K$ be an arithmetic
variety of relative dimension $e_0$, $\overline L$ a Hermitian line bundle on $X$ and
$s_1,\dots,s_m\in\Gamma(X,L)$ sections which embed $X$, such that for every $\sigma$ the weight of
$\overline L$ at $\sigma$ is the Fubini--Study weight of $(s,M_\sigma)$. Let $f$ be a Chow form of the
generic ideal $I_X$ in $e_0+1$ blocks which is multihomogeneous of degree $\delta$. If $\overline L$ has
self-intersection $a$, then
\[
a=[K:\Q]\,\covolHch^H(e_0+1,\delta,f).
\]
\formallabel{ChowFormHeightComparison₂.CF.stMain_of}\formallabel{ChowFormHeightComparison₂.CF.StMain}
\end{theorem}

\begin{proof}
Let $h$ be an arithmetic height function with $h(\eta_X)=a$. We prove by induction on $e\ge0$ the
following assertion $P(e)$: \emph{for every integral closed subscheme $Z\subseteq X$ of dimension $e+1$
with $\covolPt(\eta_Z)\cap\OO_K=0$, every $\delta_Z$ and every Chow form $f_Z$ of
$I_Z=\covolPt(\eta_Z)\otimes K$ in $e+1$ blocks which is multihomogeneous of degree $\delta_Z$, we have
$h(\eta_Z)=[K:\Q]\,\covolHch^H(e+1,\delta_Z,f_Z)$.} The theorem is $P(e_0)$ for $Z=X$; here
$\covolPt(\eta_X)\cap\OO_K=0$ because $X$ is flat over $\Z$.

Fix $e$, $Z$, $\delta_Z$ and $f_Z$ as in $P(e)$, and assume $P(e-1)$ if $e\ge1$. By
\cref{pf:covol:lem:emb}(b) the restrictions of the $s_k$ embed $Z$ with the same point ideals; since
$\covolPt(\eta_Z)\cap\OO_K=0$, $\Gamma(Z,\OO_Z)$ has characteristic $0$, $Z$ is an arithmetic variety of
relative dimension $e$, and $\covolPt(\eta_Z)=I_Z\cap\OO_K[X]$, so $I_Z$ is a projective prime. As the
$s_k$ have no common zero, there is $j$ with $\eta_Z\in D(s_j)$; put $u=(0,\dots,1,\dots,0)$, the $j$-th
unit vector, so that $s_u=s_j$, $\ell_u=X_j\notin\covolPt(\eta_Z)$ (\cref{pf:covol:lem:emb}(a)) and hence
$\ell_u\notin I_Z$.

\emph{Step 1: the arithmetic recursion.} We claim
\begin{equation}\label{pf:covol:eq:step}
h(\eta_Z)=\sum_z\ord_z(s_j|_Z)\,h(z)-\sum_{\sigma\in\Sigma_K}\covolCyc_{M_\sigma}\bigl([\sigma I_Z],e,
\covolFs_{M_\sigma}(\sigma u;\cdot)\bigr).
\end{equation}
If $e\ge1$, this is axiom (R) for $Z$, $m=1$ and $r=s_j|_Z$, since on $Z_\sigma(\C)$ the function
$\log\|s_j\|$ is $\covolFs_{M_\sigma}(\sigma u;s(y))$ (Fubini--Study weights) and its integral against
$c_1(\overline L|_Z)^e$ is the cone integral by \cref{pf:covol:prop:arch} (the restriction of $h$ to
the subschemes of $Z$ is an arithmetic height function for $\overline L|_Z$, directly from the
definition). If $e=0$, $Z$ is an arithmetic curve; every complex
point of $Z$ lies over $\eta_Z$, hence in $D(s_j)$, and axiom (P) with $r=s_j|_Z$ gives
$h(\eta_Z)=\sum_z\ord_z(s_j)\log\#\kappa(z)-\sum_\sigma\sum_y\log\|s_j(y)\|$. Every $z$ with
$\ord_z(s_j)\ne0$ is a closed point with finite residue field, so $h(z)=\log\#\kappa(z)$ by axiom (V)
with $d=0$, and the archimedean sum is the cone integral by the case $e=0$ of
\cref{pf:covol:prop:arch}.

\emph{Step 2: the components of the divisor.} By \cref{pf:covol:lem:emb}(a),(c), the points $z$ with
$\ord_z(s_j|_Z)\ne0$ correspond bijectively, via $z\mapsto\covolPt(z)$, to the relevant minimal primes of
$J=\covolPt(\eta_Z)+(\ell_u)$, and $\ord_z(s_j|_Z)=\covolLm(J,\covolPt(z))$. In particular there are
finitely many such $z$, and each has codimension one in $Z$, hence dimension $e$ by
\cref{pf:covol:lem:emb}(d). Such a $z$ is \emph{vertical} if $\covolPt(z)\cap\OO_K$ is a maximal ideal
$\mathfrak q$ and \emph{horizontal} if $\covolPt(z)\cap\OO_K=0$.

\emph{Step 3: vertical components.} If $z$ is vertical over $\mathfrak q$, then axiom (V) and
\cref{pf:covol:prop:vert}, applied to the closure of $z$, give $h(z)=\covolCdeg(\covolPt(z)\bmod
\mathfrak q,e+1)\log\Nm(\mathfrak q)$. Since $I_Z\cap\OO_K[X]=\covolPt(\eta_Z)$, the vertical components
of Step 2 over $\mathfrak q$ are exactly the elements of $\mathcal V_{\mathfrak q}$ in the definition of
$\covolVt$, with the same multiplicities, and summing over $\mathfrak q$ gives
\[
\sum_{z\text{ vertical}}\ord_z(s_j|_Z)\,h(z)=\covolVt(I_Z,e+1,u).
\]

\emph{Step 4: horizontal components.} By \cref{pf:covol:lem:chowcalc}(a), $f_Z(u,\cdot)=c\prod_W
g_W^{m_W}$, where $c\in K^\times$, $W$ runs over the support of the cycle $[I_Z+(\ell_u)]$ with
multiplicities $m_W$, and $g_W$ is a Chow form of $W$ in $e$ blocks. The primes $W$ are the ideals
$\covolPt(z)\otimes K$ for the horizontal $z$ of Step 2, and $m_W=\ord_z(s_j|_Z)$ (localisation at
$\covolPt(z)$ commutes with inverting the nonzero integers). If $e=0$, a horizontal $z$ would give a
projective prime $W$ with a Chow form in $0$ blocks, which is impossible by
\cref{pf:covol:lem:chowcalc}(d); so there is no horizontal $z$, $f_Z(u,\cdot)=c$ and
$\covolHch^H(0,\delta_Z,c)=0$ by \cref{pf:covol:lem:chowcalc}(f). If $e\ge1$, each $g_W$ divides
$f_Z(u,\cdot)$, which is multihomogeneous of degree $\delta_Z$, so $g_W$ is multihomogeneous of some
degree $\delta_W$ (\cref{pf:covol:lem:chowcalc}(e)), and comparing degrees, $\delta_Z=\sum_Wm_W\delta_W$.
The induction hypothesis $P(e-1)$, applied to the closure of $z$ (of dimension $e$), gives
$h(z)=[K:\Q]\covolHch^H(e,\delta_W,g_W)$, and \cref{pf:covol:lem:chowcalc}(f) gives
$\covolHch^H(e,\delta_Z,f_Z(u,\cdot))=\sum_Wm_W\covolHch^H(e,\delta_W,g_W)$. In both cases
\[
\sum_{z\text{ horizontal}}\ord_z(s_j|_Z)\,h(z)=[K:\Q]\,\covolHch^H\bigl(e,\delta_Z,f_Z(u,\cdot)\bigr).
\]

\emph{Step 5: conclusion.} Inserting Steps 3 and 4 in \eqref{pf:covol:eq:step} and comparing with
\cref{pf:covol:prop:recursion} (for $I=I_Z$, $f=f_Z$, $\delta=\delta_Z$) gives
$h(\eta_Z)=[K:\Q]\covolHch^H(e+1,\delta_Z,f_Z)$, which is $P(e)$.
\end{proof}

To pass from the coordinates of \cref{pf:covol:thm:main} to those of \cref{pf:covol:def:formht} we use two
further facts.

\begin{lemma}[Closures of points; change of frame]\label{pf:covol:lem:modelform}
\begin{enumerate}[label=(\alph*)]
\item Let $\Omega\supseteq K$ be a field and $t\in\Omega^m$. Let $Z$ be an integral scheme, separated,
  quasi-compact and locally of finite type over $\OO_K$, with a point $\Spec\Omega\to Z$ onto the generic
  point identifying the function field with $\Omega$ (a \emph{rational model} of $\Omega$), which is
  the \emph{closure of the point $t$}: every point of $Z$ has an index $i$ with $t_i\ne0$ and all
  $t_j/t_i$ regular at it, and for every $i$ with $t_i\ne0$ the set $U_i$ of such points is an affine
  open subset with $\Gamma(U_i,\OO_Z)=\OO_K[t_j/t_i:j\le m]$. Let $\mathcal L$ be a line bundle on
  $Z$ with a generic trivialisation and sections $s_1,\dots,s_m$ without common zero whose values in the
  generic trivialisation are $t_1,\dots,t_m$, and assume $\Gamma(Z,\OO_Z)$ has characteristic $0$. Then
  the $s_k$ embed $Z$, and $I_Z$ is the ideal of polynomials $P\in K[X_1,\dots,X_m]$ with $P(t)=0$ (the
  \emph{cone ideal} of $t$).
\item Let $y\in\Omega^N$, $g\in K^{N\times m}$ whose columns span $K^N$, $A_\sigma\in\GL_N(\C)$, and let
  $\Phi$ be an integral Chow form of the cone ideal of $y$ of dimension $d$ and degree $\delta$. Let
  $H=(M_\sigma)$ be archimedean data on $\C^m$ with $|g^{\mathsf T}x|^*_{M_\sigma}=
  |A_\sigma^{\mathsf T}x|$ for all $x\in\C^N$, where $(g^{\mathsf T}x)_k=\sum_j\sigma(g_{jk})x_j$.
  Then $\Phi\circ g$ is a Chow form in $d+1$ blocks of the cone ideal of $g^{\mathsf T}y$, it is
  multihomogeneous of degree $\delta$, and $\covolHform(\delta,g,A,\Phi)=
  \covolHch^H(d+1,\delta,\Phi\circ g)$.
\end{enumerate}
\formallabel{ChowFormHeightComparison₂.CF.stModel}\formallabel{ChowFormHeightComparison₂.CF.stForm}%
\formallabel{PositivityTheory.UnitarySelfGeneration.UnitProductCombination.formHeight_gen}
\end{lemma}

\begin{proof}
(a) On $D(s_k)$ the ratios $s_j/s_k$ are the functions $t_j/t_k$, which generate the coordinate ring by
the closure condition; and a polynomial lies in $\covolPt(\eta_Z)\otimes K$ if and only if it vanishes at
$t$, by definition of the point ideal at the generic point.

(b) Choose $P\in K^{m\times N}$ with $gP=1$ and put $z=g^{\mathsf T}y$. Since $g^{\mathsf T}P^{\mathsf T}z=z$,
the linear forms expressing $x'=g^{\mathsf T}P^{\mathsf T}x'$ lie in the cone ideal of $z$, so
$x\mapsto g^{\mathsf T}x$ and $x'\mapsto P^{\mathsf T}x'$ are inverse bijections between the zeros of the
cone ideals of $y$ and of $z$, and under them the incidence conditions $(gu_i)\cdot x=0$ and
$u_i\cdot x'=0$ correspond. Hence the zero set of $\Phi\circ g$ is the incidence set of the cone ideal of
$z$; $\Phi\circ g$ is irreducible because $(\Phi\circ g)\circ P=\Phi$, and it is multihomogeneous of
degree $\delta$. So $\Phi\circ g$ is an integral Chow form, hence a Chow form. The Gauss norms of both
sides are those of $\Phi\circ g$. For the archimedean terms fix
$\sigma$, write $M=M_\sigma$ and let $B=(M^{\mathsf T})^{-1}g^{\mathsf T}(A_\sigma^{\mathsf T})^{-1}
\colon\C^N\to\C^m$; the hypothesis says that $B$ is an isometric embedding, and
$\sigma(\Phi\circ g)(M^{-1}v)=\sigma\Phi(A_\sigma B^{\mathsf T}v)$. If $v$ is uniformly distributed on
the unit sphere of $\C^m$, then $\pi(v)=B^{\mathsf T}v$ is the orthogonal projection onto an
$N$-dimensional subspace (up to complex conjugation of coordinates), $\pi(v)/|\pi(v)|$ is uniformly
distributed on the unit sphere of $\C^N$ and independent of $|\pi(v)|^2$, which has the beta
distribution with parameters $(N,m-N)$, so that $\mathbb E\log|\pi(v)|^2=\psi(N)-\psi(m)=
\covolHar_{N-1}-\covolHar_{m-1}$ ($\psi$ the digamma function). By homogeneity of degree $\delta$ in each
of the $d+1$ blocks,
\[
\covolMah_M\bigl(\sigma(\Phi\circ g)\bigr)=\covolMah^{A_\sigma}(\sigma\Phi)+
\frac{(d+1)\delta}2\bigl(\covolHar_{N-1}-\covolHar_{m-1}\bigr),
\]
which is the asserted identity after adding $\frac{(d+1)\delta\covolHar_{m-1}}2$.
\end{proof}

\begin{theorem}[Chow height formula]\label{pf:covol:thm:chowformheight}
Assume the Hodge package $\mathcal P$. Let $\Omega\supseteq K$ be a field,
$d,N,m,\delta\in\N$, $g\in K^{N\times m}$ with columns spanning $K^N$, and $A_\sigma\in\GL_N(\C)$
$(\sigma\in\Sigma_K)$ conjugation compatible. Let $y\in\Omega^N$ and put $z_k=\sum_jg_{jk}y_j$. Let $Z$ be
a rational model of $\Omega$ over $\OO_K$ which is the closure of $z$ and an arithmetic variety of
relative dimension $d$; let $\Phi$ be an integral Chow form of the cone ideal of $y$ of dimension $d$ and
degree $\delta$; let $\mathcal L$ be a line bundle on $Z$ with a generic trivialisation and sections
$s_1,\dots,s_m$ without common zero whose values are $z_1,\dots,z_m$; and let $M_\sigma\in\GL_m(\C)$
$(\sigma\colon\OO_K\to\C)$ with $M_{\bar\sigma}=\overline{M_\sigma}U_\sigma$ for unitary $U_\sigma$ and
\[
\sum_a\Bigl|\sum_kM_\sigma(k,a)\sum_j\sigma(g_{jk})x_j\Bigr|^2=|A_\sigma^{\mathsf T}x|^2\qquad
(x\in\C^N).
\]
Let $\overline{\mathcal L}_{\mathrm{FS}}$ be $\mathcal L$ with the weights $\log\sum_a|\sum_kM_\sigma(k,a)s_k|^2$.
Then $\overline{\mathcal L}_{\mathrm{FS}}$ has self-intersection $[K:\Q]\,\covolHform(\delta,g,A,\Phi)$.
\formallabel{ChowFormHeightComparison₂.CF.glue}\formallabel{ChowFormHeightComparison₂.CF.chowFormHeight_of_arch}%
\formallabel{ChowFormHeightComparison₂.chow_form_height_bound}%
\formallabel{PositivityTheory.UnitarySelfGeneration.ChowFormHeight}
\end{theorem}

\begin{proof}
By \cref{pf:covol:lem:modelform}(a) (with $t=z$; $\Gamma(Z,\OO_Z)$ has characteristic $0$ because $Z$ is
flat over $\Z$) the $s_k$ embed $Z$ and $I_Z$ is the cone ideal of $z$. Let $H$ be the archimedean data
given by the matrices $M_\sigma^{-1}$, so that $|v|^*_{M_\sigma^{-1}}=(\sum_a|\sum_kM_\sigma(k,a)v_k|^2)^{1/2}$;
the relation $M_{\bar\sigma}=\overline{M_\sigma}U_\sigma$ with $U_\sigma$ unitary gives the conjugation
condition. By construction the weight of $\overline{\mathcal L}_{\mathrm{FS}}$ at $\sigma$ is the
Fubini--Study weight of $(s,M_\sigma^{-1})$, and the displayed hypothesis says
$|g^{\mathsf T}x|^*_{M_\sigma^{-1}}=|A_\sigma^{\mathsf T}x|$. By \cref{pf:covol:lem:modelform}(b),
$\Phi\circ g$ is a Chow form of $I_Z$ in $d+1$ blocks, multihomogeneous of degree $\delta$, and
$\covolHform(\delta,g,A,\Phi)=\covolHch^H(d+1,\delta,\Phi\circ g)$. By the existence of arithmetic height
functions (part of $\mathcal P$) there is an arithmetic height function $h_0$ for
$\overline{\mathcal L}_{\mathrm{FS}}$, so $\overline{\mathcal L}_{\mathrm{FS}}$ has self-intersection
$h_0(\eta_Z)$, and \cref{pf:covol:thm:main} gives $h_0(\eta_Z)=[K:\Q]\covolHch^H(d+1,\delta,\Phi\circ g)=
[K:\Q]\covolHform(\delta,g,A,\Phi)$.
\end{proof}

\begin{remark}\label{pf:covol:rem:classical}
(1) The conjugation compatibility of the frames $A_\sigma$ is a hypothesis of the formal statement which
the proof does not use.
(2) \Cref{pf:covol:thm:chowformheight} is the comparison theorem of Philippon and Soul\'e
\cite{Philippon1991,Soule1991} between the height of a projective variety for the Fubini--Study metric
and the height of its Chow form defined with Mahler measures on spheres, in the form
$h_{\overline{\OO(1)}}(X)=\sum_{\mathfrak q}\log|F_X|_{\mathfrak q}+\sum_\sigma\bigl(\covolMah(\sigma F_X)+
\frac{(d+1)\deg X}2\covolHar_{N-1}\bigr)$ for $X\subseteq\mathbb P^{N-1}$ of dimension $d$; compare
\cite[\S4.3]{BostGilletSoule1994}. For instance, for $X=\mathbb P^{N-1}$ (with $d=N-1$, $\delta=1$ and Chow form
the determinant) the right-hand side is $\frac12\sum_{k=1}^{N-1}\covolHar_k$, the classical value of the
Fubini--Study height of projective space. In the formal statement, self-intersections are those of
\cref{pf:covol:def:height}, the variety is given as a rational model which is the closure of a point, and
the metric is pulled back by a linear map $g$; the published form of the formula for subvarieties of
projective space is recalled in \cref{pf:covol:lit:fs}.
\end{remark}

% ---------------------------------------------------------------------------
\subsection{The Chow form height comparison}\label{sec:pf:covol:comparison}

\begin{definition}[Frames, image degree and image height]\label{pf:covol:def:image}
Let $V$ be a finite-dimensional $K$-vector space, $W\subseteq V$ a finitely generated $\OO_K$-submodule
and $\nu=(\nu_\sigma)_{\sigma\in\Sigma_K}$ a family of norms on $V$. Let $\Omega\supseteq K$ be a field and
$\mathrm{ev}\colon V\to\Omega$ a $K$-linear map.
\begin{enumerate}[label=(\alph*)]
\item A \emph{frame} of $(\nu,W)$ is a triple $(b,w,A)$ consisting of a basis $b_1,\dots,b_N$ of $V$, a
  matrix $w\in K^{N\times n}$ whose columns span $K^N$ and such that the vectors $\sum_jw_{jk}b_j$
  ($k\le n$) generate $W$ over $\OO_K$, and conjugation compatible invertible matrices $A_\sigma$ with
  $\nu_\sigma(x)=|A_\sigma^{-1}(\sigma x_j)_j|$, where $x=\sum_jx_jb_j$.
\item $\mathrm{ev}$ \emph{generates} $\Omega$ if $\mathrm{ev}\ne0$ and $\Omega=K(\mathrm{ev}(v)/\mathrm{ev}(v'):
  v,v'\in V)$. The \emph{image ideal} of $\mathrm{ev}$ for a basis $b$ is the cone ideal of
  $(\mathrm{ev}(b_j))_j$.
\item The \emph{image degree} $\covolImDeg(\mathrm{ev},D)$ is an integer $\delta$ such that for some basis
  $b$ the image ideal has an integral Chow form of dimension $D$ and degree $\delta$, and $0$ if there is
  no such basis. The \emph{image height} $\covolImHt(\nu,W,\mathrm{ev},D)$ is a real number $h$ such
  that $(D+1)\delta h=\covolHform(\delta,w,A,f)$ for some frame $(b,w,A)$ of $(\nu,W)$ and some integral
  Chow form $f$ of the image ideal for $b$ of dimension $D$ and degree $\delta$, and $0$ if there is no
  such $h$.
\end{enumerate}
\formallabel{GlobalTheory.CompHt.imageDeg}\formallabel{GlobalTheory.CompHt.imageHeight}
\end{definition}

Both quantities are well defined, i.e.\ independent of the choices: the degree of an integral Chow form
is the degree of the image variety, and the image height does not depend on the frame and the Chow form
(\lref{GlobalTheory.ImageHeightUnique.isImageHeight_unique}; under a change of basis $b'$ the form
height transforms by \lref{GlobalTheory.ImageHeightUnique.formHeight_linSubst}, which states
$\covolHform(\delta,w',A',f\circ P)=\covolHform(\delta,w,A,f)$ for two frames and the change-of-basis
matrix $P$ from $b$ to $b'$). If $\covolImDeg>0$, the image variety (the closure of the point
$(\mathrm{ev}(b_j))_j$ in $\mathbb P^{N-1}_K$) has dimension $D$ and degree $\covolImDeg$, and $\covolImHt$ is
its height for the Fubini--Study metrics of the frame, normalised by $(D+1)\covolImDeg$
(\cref{pf:covol:rem:classical}).

\begin{definition}[Models of images]\label{pf:covol:def:model}
In the situation of \cref{pf:covol:def:image}, let $D\in\N$, let $p\colon X\to\Spec\OO_K$ with $X$
integral, $\overline L$ a Hermitian line bundle on $X$, $\psi\colon\Omega\to\kappa(X)$ a homomorphism to
the function field and $s\colon W\to\Gamma(X,L)$ a map. Then $(X,\overline L,\psi,s)$ is a \emph{model}
of $(\nu,W,\mathrm{ev},D)$ if: $X$ is an arithmetic variety of relative dimension $D$; $X$ is normal;
$\psi$ is bijective and compatible with $\OO_K\to\Gamma(X,\OO_X)$; $\psi(\mathrm{ev}\,w)\,s(w')=
\psi(\mathrm{ev}\,w')\,s(w)$ as rational sections for all $w,w'\in W$; the sections $s(w)$ have no
common zero; and there are a frame $(b,w,A)$ of $(\nu,W)$, with generators $\gamma_k=\sum_jw_{jk}b_j$ of
$W$, and $c\in K^{n\times N}$ with $b_j=\sum_kc_{kj}\gamma_k$, such that the weight of $\overline L$ at
$\sigma$ is $\log\sum_a|\sum_k(\sum_l\sigma(c_{kl})A_\sigma(l,a))\,s(\gamma_k)|^2$.
\formallabel{GlobalTheory.ImageModel.IsModelOf}
\end{definition}

The last condition says that the metric of $\overline L$ is the Fubini--Study metric attached to the
sections $s(w)$ and to the norms $\nu$: by the ratio condition, $\sum_kc_{kl}\,s(\gamma_k)$ is the
section corresponding to $b_l$, so the weight is $\log|(s(b_l))_l|^{*2}$ for the dual norm
$|y|^*=|A_\sigma^{\mathsf T}y|$ of the frame.

\begin{lemma}[Scaling]\label{pf:covol:lem:scaling}
If $f\ne0$ is multihomogeneous of degree $\delta$ in $B$ blocks of $N\ge1$ variables and $A_\sigma$ are
invertible, then $\covolHform(\delta,w,e^tA,f)=\covolHform(\delta,w,A,f)+B\delta t$ for $t\in\R$.
\formallabel{PositivityTheory.UnitarySelfGeneration.FrameHeightScaling.formHeight_smul}
\end{lemma}

\begin{proof}
$\sigma f(e^tA_\sigma z)=e^{B\delta t}\sigma f(A_\sigma z)$, so each of the $[K:\Q]$ archimedean terms
increases by $B\delta t$, and the sum is divided by $[K:\Q]$.
\end{proof}

\begin{proposition}[Self-intersection of the Fubini--Study model]\label{pf:covol:prop:vol}
Assume the Hodge package $\mathcal P$. In the situation of \cref{pf:covol:def:image}, let $(b,g,A)$ be a
frame of $(\nu,W)$ with $g\in K^{N\times m}$, let $\mathrm{ev}$ generate $\Omega$ and
$\covolImDeg(\mathrm{ev},d)>0$. Let $Y$ be a rational model of $\Omega$ over $\OO_K$ which is an arithmetic
variety of relative dimension $d$, such that for $z_k=\mathrm{ev}(\sum_jg_{jk}b_j)$ every point of $Y$ has
an index $i$ with $z_i\ne0$ and all $z_j/z_i$ regular at the point; let $\mathcal L$ be a line bundle on
$Y$ with a generic trivialisation and sections $s_1,\dots,s_m$ without common zero with values $z_k$; and
let $M_\sigma$ be as in \cref{pf:covol:thm:chowformheight}, with
$\sum_a|\sum_kM_\sigma(k,a)\sum_j\sigma(g_{jk})x_j|^2=|A_\sigma^{\mathsf T}x|^2$. Then for every $t\in\R$
the Hermitian line bundle $\overline{\mathcal L}_{\mathrm{FS}}(t)$ has self-intersection
\[
(d+1)\,[K:\Q]\,\covolImDeg(\mathrm{ev},d)\,\bigl(\covolImHt(\nu,W,\mathrm{ev},d)+t\bigr).
\]
\formallabel{PositivityTheory.UnitarySelfGeneration.volume_le_self_intersection}
\end{proposition}

\begin{proof}
The values $z_k$ generate $\Omega$ over $K$ by their ratios, because $\mathrm{ev}$ does and the columns of
$g$ span $K^N$. By \lref{GlobalTheory.ZH.exists_isCl_isProjOver}, which states that for a
Noetherian ring $o\subseteq\Omega$ and $t\in\Omega^{n+1}$ with some $t_i\ne0$ such that $o$ and the ratios
$t_j/t_i$ generate the field $\Omega$, there is a rational model of $\Omega$ over $o$ which is the closure
of the point $t$ (\cref{pf:covol:lem:modelform}(a)) and is projective over $o$, there is such a model $Z$
for $o=\OO_K$ and $t=z$. Its ring of functions has characteristic $0$, so it is flat over $\Z$; and it has
dimension $d+1$, because the transcendence degree of $\Omega$ over $K$ is $d$ when the image ideal has an
integral Chow form of dimension $d$. So $Z$ is an arithmetic variety of relative dimension $d$. Since
every point of $Y$ has a coordinate $z_i$ with regular ratios $z_j/z_i$, the identity of $\Omega$ induces
a morphism $\pi\colon Y\to Z$ over $\OO_K$ which maps the generic point to the generic point with residue
degree $1$. The coordinate functions give a line bundle $\mathcal L'$ on $Z$ with a generic trivialisation
and sections $s'_k$ with values $z_k$.

Let $\delta=\covolImDeg(\mathrm{ev},d)$ and choose an integral Chow form $f'$ of dimension $d$ and degree
$\delta$ of the image ideal for some basis $b'$; transporting it by the change of basis gives an integral
Chow form $f$ of the image ideal for $b$, i.e.\ of the cone ideal of $y=(\mathrm{ev}(b_j))_j$. By
definition of the image height and its independence of choices,
$\covolImHt(\nu,W,\mathrm{ev},d)=\covolHform(\delta,g,A,f)/((d+1)\delta)$.

Apply \cref{pf:covol:thm:chowformheight} to $Z$, the sections $s'_k$, the frames $e^tA_\sigma$ and the
matrices $e^tM_\sigma$ (the hypotheses are homogeneous in this scaling): the Fubini--Study bundle
$\overline{\mathcal L'}_{\mathrm{FS}}$ of $(s',e^tM)$ has self-intersection
$[K:\Q]\covolHform(\delta,g,e^tA,f)$. By the projection formula of $\mathcal P$, applied to $\pi$ with
residue degree $1$, $\pi^*\overline{\mathcal L'}_{\mathrm{FS}}$ has the same self-intersection. The
sections $\pi^*s'_k$ and $s_k$ have the same values, and multiplying $M_\sigma$ by $e^t$ adds $2t$ to the
Fubini--Study weights; so the identity is an isometry from $\pi^*\overline{\mathcal L'}_{\mathrm{FS}}$ to
$\overline{\mathcal L}_{\mathrm{FS}}(t)$, and by \cref{pf:covol:lem:unique} the latter has the same
self-intersection. Finally, by \cref{pf:covol:lem:scaling} with $B=d+1$,
\[
[K:\Q]\covolHform(\delta,g,e^tA,f)=[K:\Q]\bigl(\covolHform(\delta,g,A,f)+(d+1)\delta t\bigr)=
(d+1)[K:\Q]\delta\bigl(\covolImHt+t\bigr).\qedhere
\]
\end{proof}

\begin{theorem}[Chow form height comparison]\label{pf:covol:thm:comparison}
Assume the Hodge package $\mathcal P$. Let $K$ be a totally complex number field, $V$ a $K$-vector space,
$W\subseteq V$ a finitely generated $\OO_K$-submodule with $KW=V$, $\nu$ a family of norms on $V$ which
has a Gram frame (\cref{def:covolume:slopes}) and satisfies $\nu_{\bar\sigma}=\nu_\sigma$,
$\mathrm{ev}\colon V\to\Omega$ a $K$-linear map which generates $\Omega$, and $D\in\N$ with
$\covolImDeg(\mathrm{ev},D)>0$. Let $(X,\overline L,\psi,s)$ be a model of $(\nu,W,\mathrm{ev},D)$, let
$t\in\R$, and suppose that $\overline L(t)$ has self-intersection $a$. Then
\[
a\le[K:\Q]\,(D+1)\,\covolImDeg(\mathrm{ev},D)\,\bigl(\covolImHt(\nu,W,\mathrm{ev},D)+t\bigr).
\]
In fact equality holds.
\formallabel{ChowFormHeightComparison₂.chow_form_height_comparison}%
\formallabel{PositivityTheory.UnitarySelfGeneration.weil_height_inequality}
\end{theorem}

\begin{proof}
Since $\psi$ is bijective and compatible with $\OO_K$, $X$ with the point defined by $\psi^{-1}$ is a
rational model of $\Omega$ over $\OO_K$. Let $(b,g,A)$, $\gamma_k$ and $c$ be the frame data of
\cref{pf:covol:def:model}. There are matrices $M_\sigma\in\GL_m(\C)$ with
$M_{\bar\sigma}=\overline{M_\sigma}U_\sigma$, $U_\sigma$ unitary, and
$\sum_a|\sum_kM_\sigma(k,a)\sum_j\sigma(g_{jk})x_j|^2=|A_\sigma^{\mathsf T}x|^2$: this is linear algebra,
since the map $x\mapsto(\sum_j\sigma(g_{jk})x_j)_k$ is injective and the Hermitian form
$|A_\sigma^{\mathsf T}x|^2$ on its image extends to a positive definite form on $\C^m$, compatibly with
complex conjugation. By the ratio condition and the absence of common zeros, the generic point of $X$ lies
in the non-vanishing locus of some $s(w)$, and $L$ has a generic trivialisation in which every $s(w)$ has
value $\mathrm{ev}(w)$; in particular every point of $X$ has a coordinate $\mathrm{ev}(\gamma_i)\ne0$ with
all ratios regular. \Cref{pf:covol:prop:vol}, applied to $Y=X$, $\mathcal L=L$ and the sections
$s(\gamma_k)$, shows that the Fubini--Study bundle $\overline L'_{\mathrm{FS}}(t)$ of $(s(\gamma_k),M)$ has
self-intersection $(D+1)[K:\Q]\covolImDeg(\covolImHt+t)$. The metric of $\overline L$ in the definition of
a model is the Fubini--Study metric of the same sections for the dual norm $|A_\sigma^{\mathsf T}\cdot|$
(by the ratio condition, $\sum_kc_{kl}s(\gamma_k)$ is the section corresponding to $b_l$), so the
identity is an isometry between $\overline L(t)$ and $\overline L'_{\mathrm{FS}}(t)$. By
\cref{pf:covol:lem:unique}, $a=(D+1)[K:\Q]\covolImDeg(\covolImHt+t)$.
\end{proof}

\begin{remark}\label{pf:covol:rem:equality}
The formal statement is an inequality, but its proof establishes the equality. The hypotheses that $K$
is totally complex, that $\nu$ has a Gram frame and is conjugation invariant, and that $W$ is finitely
generated are part of the statement but are not used in the proof; the frame data needed are contained in
the definition of a model. The theorem is the self-intersection hypothesis of the lift estimates
(\cref{pf:covol:def:packages}).
\end{remark}


% ---------------------------------------------------------------------------
\subsection{Small forms and small sections}\label{sec:pf:covol:small}

We use the deep covers of \cref{sec:covolume:lift}.

\begin{definition}[Small forms; admissible components]\label{pf:covol:def:small}
Fix a global datum $\mathcal D$ with host field $E$ (a CM field), a core datum $T$ of $\mathcal D$, an
integer $j$ and a deep cover $Y$ of $T$ at level $j$ which is \emph{good} for some parameters (the
condition $\mathrm{GoodCov}$ of \cref{sec:covolume:lift}). The cover carries a unitary moduli space $X^o$
over a number field $F$ (its coefficient field) with complex uniformisations at the good embeddings
$\sigma\colon F\to\C$. Automorphic forms are \emph{rules}: functions which assign to every framed point of
$X^o$ with values in an $F$-algebra $S$ (an $S$-point with a frame of the Lie algebra of the universal
abelian scheme and a generator of the polarisation line) an element of $S$, compatibly with base change.
For an irreducible component $Z$ of the generic fibre of $X^o$ and a weight $G$:
\begin{enumerate}[label=(\alph*)]
\item $\mathcal C(Z,G)$ is the $F$-vector space of cusp forms of weight $G$ supported on $Z$,
  $\Lambda(Z,G)\subseteq\mathcal C(Z,G)$ its $\OO_F$-lattice of integral forms, and $\nu_\sigma(Z,G)$ the
  Petersson-type norm at $\sigma$ (computed with a probability measure on the frames over $Z$);
\item $\Omega_Z$ is the field of rational functions in the flag variables over the function field of $Z$,
  and $\covolGen_Z$ is the $F$-algebra homomorphism from automorphic forms to $\Omega_Z$ given by evaluation
  at a generic framed point over $Z$; its restriction to $\mathcal C(Z,G)$ is the \emph{evaluation map} of
  the lattice;
\item for $\tau\in\N$ and $t\in\R$, the set $\covolSmall(Z,G,\tau,t)$ of \emph{small forms} consists of the
  $f\in\mathcal C(Z,\tau G)$ such that $v(f)\ge\tau\inf_{w\in\Lambda(Z,G)}v(w)$ for every maximal ideal
  $\mathfrak l$ of $\OO_F$ and every valuation $v$ of the family of valuations attached to $\mathfrak l$,
  and such that at every good embedding $\sigma$ and every frame point $y$ over $Z$ the pointwise norm of
  $f$ satisfies $|f(y)|^2\le e^{2\tau t}B_G(y)^\tau$, where $B_G$ is the Bergman function of weight $G$ for
  the Petersson norm on $Z$.
\end{enumerate}
For a weight $u$ of the datum of $T$ at level $j$ let $G_u$ be the associated weight of the deep cover
(with shift $0$), and let $\deg(Y,u,Z)=\covolImDeg(\covolGen_Z|_{\mathcal C(Z,G_u)},\dim\mathrm{Fl})$ be the
image degree of the evaluation map in the dimension $\dim\mathrm{Fl}$ of the flag variety of the datum. The
pair $(Z,u)$ is \emph{admissible} if $(\nu(Z,G_u),\Lambda(Z,G_u))$ has a Gram frame, $\Lambda(Z,G_u)$ is
finitely generated and spans $\mathcal C(Z,G_u)$, the evaluation map generates $\Omega_Z$, the lattice has
no base locus along the valuations of these families (for every such valuation $v$ and every form $f$ of
any weight with $v(f)\ne\infty$ there is $w\in\Lambda(Z,G_u)$ with $v(w)\ne\infty$), and $\deg(Y,u,Z)>0$.
\formal{GlobalTheory.CovTop.Deep.Good}\formal{GlobalTheory.AllWt.cuspOn}\formal{GlobalTheory.AllWt.latOn}%
\formal{GlobalTheory.AllWt.normOn}\formal{GlobalTheory.CompHt.genVal}\formal{GlobalTheory.CovHt.smallL}%
\formal{GlobalTheory.AllWt.wtOf}\formallabel{GlobalTheory.AllWt.GoodOn}\formallabel{GlobalTheory.DegVal.BirOn}%
\formallabel{GlobalTheory.CovHt.NoBase}
\end{definition}

\begin{theorem}[Small forms give small sections]\label{pf:covol:thm:small}
Assume the non-vanishing statement (NV) of \cref{pf:covol:thm:nv} and the Hodge package $\mathcal P$. Let
$Y$ be a good deep cover as in \cref{pf:covol:def:small}, with $[E:\Q]/2>2$, let $(Z,u)$ be admissible,
and let $(X,\overline L,\psi,s)$ be a model (\cref{pf:covol:def:model}) of
$(\nu(Z,G_u),\Lambda(Z,G_u),\covolGen_Z,\dim\mathrm{Fl})$. Then for every $t\in\R$ and every $i\ge1$,
\[
\#\covolSmall(Z,G_u,8i,t)\le\#\bigl\{\text{small sections of }\overline L(t)^{\otimes8i}\bigr\}
\]
(cardinalities in $\N\cup\{\infty\}$).
\formallabel{PositivityTheory.Xj7N.small_sections_injective}
\end{theorem}

The theorem, applied with (NV), is the small-sections hypothesis of the lift estimates
(\cref{pf:covol:def:packages}). Its proof constructs an injection $f\mapsto S_f$ and needs three lemmas.

\begin{lemma}[Rational sections]\label{pf:covol:lem:ratsmall}
Let $X$ be an integral scheme over $\Spec k$, $\overline L$ a Hermitian line bundle on $X$, $s_0$ a section
of $L$ not vanishing at $\eta_X$, $\tau\in\N$ and $t\in\R$. There is an injective map $R$ from the sections
of $L^{\otimes\tau}$ to the function field $\kappa(X)$ with the following property: if $h\in\kappa(X)$ is
such that the rational section $hs_0^\tau$ is regular at every point of $X$ and
$\|hs_0^\tau(y)\|^2\le e^{2\tau t}$ (norm of $\overline L^{\otimes\tau}$) at every complex point $y$ at which
it is regular, then $h=R(S)$ for a small section $S$ of $\overline L(t)^{\otimes\tau}$.
\formallabel{PositivityTheory.Xj7N.Inj.exists_rat_small}
\end{lemma}

\begin{proof}
Let $R(S)=S/s_0^\tau$. If $hs_0^\tau$ is regular everywhere, it is a global section $S$ with $R(S)=h$,
and its norm for $\overline L(t)^{\otimes\tau}$ is $e^{-\tau t}\|hs_0^\tau\|\le1$.
\end{proof}

\begin{lemma}[From generic points to all complex points]\label{pf:covol:lem:dense}
Let $(X,\overline L,\psi,s)$ be a model of $(\nu,W,\mathrm{ev},d)$, $\beta\in\Omega$, $\tau\in\N$ and
$t\in\R$. Suppose that for every frame $(b,g,A)$ of $(\nu,W)$, every $\sigma\in\Sigma_K$ and every field
embedding $\rho\colon\Omega\to\C$ extending $\sigma$,
\begin{equation}\label{pf:covol:eq:genbound}
|\rho(\beta)|\le e^{\tau t}\,\bigl|A_\sigma^{\mathsf T}(\rho(\mathrm{ev}\,b_j))_j\bigr|^{\tau}.
\end{equation}
Let $w\in W$ with $\mathrm{ev}(w)\ne0$ and $F=\psi(\beta)\psi(\mathrm{ev}\,w)^{-\tau}s(w)^\tau$. Then
$\|F(y)\|^2\le e^{2\tau t}$ (norm of $\overline L^{\otimes\tau}$) at every complex point $y$ of $X$ at which
$F$ is regular.
\formallabel{ChowFormHeightComparison₂.GenericDenseBound.norm_le_of_emb}%
\formallabel{PositivityTheory.Xj7N.weil_height_comparison_of_degree_bound}
\end{lemma}

\begin{proof}
Let $(b,g,A)$, $\gamma_k$ and $c$ be the frame data of the model, and let $y$ lie over
$\sigma_0\colon\OO_K\to\C$. Choose an affine open set $O\ni y$ in a trivialising open set of $L$ on which
$F$ is regular, and let $B=\Gamma(O,\OO_X)$, a domain of finite type over $\Z$ with fraction field
$\kappa(X)=\psi(\Omega)$. In the trivialisation, $\|F(y)\|^2=|F(y)|^2e^{-\tau\varphi(y)}$ with
$\varphi=\log\sum_a|\sum_k(\sum_l\sigma(c_{kl})A_\sigma(l,a))\,s(\gamma_k)|^2$, and the desired inequality is
an inequality $\Phi_\sigma(r(y))\le\Psi_\sigma(r(y))$ between continuous functions of the values at $y$ of
the finitely many elements $r=(F,s(\gamma_1),\dots)$ of $B$.

\emph{The inequality holds at every embedding $\rho\colon\kappa(X)\to\C$.} Let $\sigma=\rho\psi|_K$ and
$\lambda=\rho(\psi(\mathrm{ev}\,w)^{-1}s(w))$. By the ratio condition,
$\rho(s(\gamma_k))=\rho\psi(\mathrm{ev}\,\gamma_k)\lambda$ and $\rho(F)=\rho\psi(\beta)\lambda^\tau$; since
$b_l=\sum_kc_{kl}\gamma_k$, $\sum_k(\sum_l\sigma(c_{kl})A_\sigma(l,a))\rho\psi(\mathrm{ev}\,\gamma_k)=
\sum_lA_\sigma(l,a)\rho\psi(\mathrm{ev}\,b_l)$. The inequality at $\rho$ is therefore
\eqref{pf:covol:eq:genbound} for $\rho\psi$, squared and multiplied by $|\lambda|^{2\tau}$.

\emph{Limits.} Write $B$ as a quotient of $\Z[X_1,\dots,X_n]$; the ring homomorphisms $B\to\C$ form a
closed algebraic subset $V\subseteq\C^n$ defined over $\Q$, whose irreducible components are conjugate
under automorphisms of $\C$ because $B\otimes\Q$ is a domain. For $0\ne b\in B$ the zero set of $b$ in $V$
contains no component (otherwise it would contain all of them, and $b$ would be $0$ in the domain $B$),
so it is nowhere dense in $V$. Since $B$ is countable, Baire's theorem shows that the points of $V$ at
which no nonzero element of $B$ vanishes are dense; at such a point evaluation is injective on $B$ and
extends to a field embedding of $\kappa(X)$. Hence $y$ is the pointwise limit of a sequence of field
embeddings $\rho_n\colon\kappa(X)\to\C$. As $\Sigma_K$ is finite, a subsequence has a constant restriction
$\sigma$ to $K$, and $\sigma|_{\OO_K}=\sigma_0$ by continuity. Passing to the limit in
$\Phi_\sigma(r(\rho_n))\le\Psi_\sigma(r(\rho_n))$ gives the inequality at $y$.
\end{proof}

\begin{lemma}[Regularity]\label{pf:covol:lem:regular}
Let $(X,\overline L,\psi,s)$ be a model of $(\nu,W,\mathrm{ev},d)$, $\beta\in\Omega$ and $\tau\in\N$.
Suppose that for every maximal ideal $\mathfrak l$ of $\OO_K$ and every valuation ring $O\supseteq\OO_K$ of
$\Omega$ whose maximal ideal contains $\mathfrak l$ there is $l\in W$ with $\beta\in\mathrm{ev}(l)^\tau O$.
Then for every $w\in W$ with $\mathrm{ev}(w)\ne0$ the rational section
$\psi(\beta)\psi(\mathrm{ev}\,w)^{-\tau}s(w)^\tau$ of $L^{\otimes\tau}$ is regular at every point of $X$.
\formallabel{ChowFormHeightComparison₂.RegularPointExistence.reg_of_val}%
\formallabel{ChowFormHeightComparison₂.SmallGenerators.reg_of_val}
\end{lemma}

\begin{proof}
Let $x\in X$ and $F=\psi(\beta)\psi(\mathrm{ev}\,w)^{-\tau}s(w)^\tau$, written in a trivialisation near
$x$. $X$ is normal, so the local ring at $x$ is the intersection of the valuation rings of $\kappa(X)$
containing it, and it suffices to show $F\in O$ for each such $O$. The valuation ring
$\psi^{-1}(O)$ of $\Omega$ contains $\OO_K$; it contains a valuation ring $O'\supseteq\OO_K$ of $\Omega$
whose maximal ideal contains some maximal ideal $\mathfrak l$ of $\OO_K$ (refine by a valuation of the
residue field). By hypothesis $\beta=\mathrm{ev}(l)^\tau c$ for some $l\in W$ and $c\in O'\subseteq\psi^{-1}(O)$.
By the ratio condition, $F=\psi(c)\bigl(\psi(\mathrm{ev}\,l)\psi(\mathrm{ev}\,w)^{-1}s(w)\bigr)^\tau=
\psi(c)\,s(l)^\tau$, and $s(l)$ is regular at $x$, so $F\in O$.
(The formal development contains two copies of this argument, one for the points of the generic fibre
and one for the points over the maximal ideals of $\OO_K$.)
\end{proof}

\subsubsection*{Forms with vanishing generic values}
Recall from \cref{pf:covol:def:small} that rules are functions on framed points. Formally they are
defined on the framed points with values in the \emph{test rings} $F[x_1,\dots,x_r]/\mathfrak a$ with
$\mathfrak a$ radical, and are compatible with base change. A rule $g$ has \emph{exponent}
$e\colon\mathrm{Ix}\to\Z$ if $g(y')=\prod_it_i^{e_i}\,g(y)$ whenever $y'$ is obtained from $y$ by a Borel
change of frame with diagonal $t$ (the index set $\mathrm{Ix}$ consists of the polarisation index and the
pairs $(\tau,i)$, $i<\mathrm{mult}(\tau)$).

\begin{lemma}[Rules with vanishing generic values]\label{pf:covol:lem:secz}
Let $g$ be a rule of exponent $e$ such that $\covolGen_Z(g)=0$ for every irreducible component $Z$ of the
generic fibre of $X^o$. Then $g=0$.
\formallabel{GlobalTheory.CompHt.secZ_eq_zero_of_genVal}%
\formallabel{GlobalTheory.CompHt.evR_eq_zero_of_genVal_eq_zero}
\end{lemma}

\begin{proof}
We use the \emph{atlas property} of generic frames (\lref{GlobalTheory.Reg.evLaws_genFr_atlas}): for
every field $L$ over $F$ and every framed point $y$ over $L$ there are a component $Z$, a smooth
$F$-subalgebra $R_C\subseteq\Omega_Z$ with fraction field $\Omega_Z$, a framed point $y_C$ over $R_C$ whose
image in $\Omega_Z$ is obtained from the generic framed point of $Z$ by a Borel change of frame with
diagonal $t$, an $F$-homomorphism $\psi\colon R_C\to L$, and a diagonal $t'$ such that $y$ is obtained from
$\psi_*(y_C)$ by the Borel change with diagonal $t'$.

\emph{Points over fields.} In this situation, by the transformation law and the compatibility with base
change, $g(y)=\prod_it_i'^{e_i}\,\psi(g(y_C))$, and the image of $g(y_C)$ in $\Omega_Z$ is
$\prod_it_i^{e_i}\,\covolGen_Z(g)=0$. As $R_C\subseteq\Omega_Z$, $g(y_C)=0$ and so $g(y)=0$.

\emph{Test rings.} Let $S=F[x]/\mathfrak a$ be a test ring, $y$ a framed point over $S$, and $a\in F[x]$ a
representative of $g(y)$. For a maximal ideal $\mathfrak m\supseteq\mathfrak a$, the image of $y$ under
$S\to F[x]/\mathfrak m$ is a framed point over a field (\lref{GlobalTheory.FlagTan.exists_frMap}), at
which the value of $g$ is the image of $g(y)$ (\lref{GlobalTheory.FlagTan.evR_eq}); it is $0$ by the
first step, so $a\in\mathfrak m$. Since $F[x]$ is a Jacobson ring and $\mathfrak a$ is radical,
$\mathfrak a$ is the intersection of the maximal ideals containing it; hence $a\in\mathfrak a$ and $g(y)=0$.
\end{proof}

\subsubsection*{Holomorphic sections over a component}
Fix a good embedding $\sigma$ of $F$. The complex uniformisation $V$ of $X^o$ at $\sigma$ consists of charts
$k$, each a domain of frames of a Hermitian symmetric domain, with a \emph{model weight} $\lambda_k$ on chart
$k$ attached to every weight of the datum; $\mathrm{rd}(f)(k,x)$ denotes the value of a rule $f$ at the
framed point of chart $k$ at $x$, read through the uniformisation, and the \emph{frames over $Z$} are the
pairs $(k,x)$ whose framed point lies on $Z$. A \emph{holomorphic section} of weight $\lambda$ on chart $k$ is a
holomorphic function $P$ on its frames with the transformation law of $\lambda$; it is \emph{cuspidal} if
it has moderate growth and its constant Fourier--Jacobi coefficients at all isotropic vectors vanish. For a
toroidal flag datum $\mathrm{TF}$ of $X^o$ (\lref{DifferentiableLines.nonempty_toroidal_flag_bundle},
which provides one from the first statement of the Hodge package when the flag variety has positive
dimension), with flag model $\overline{\mathrm{Fl}}$, interior $\mathrm{Fl}^o$ and automorphic line bundles
$\mathcal L(e)$, an algebraic form of exponent $e$ \emph{vanishing along the boundary} is the rule defined by
a global section of $\mathcal L(e)$ whose non-vanishing locus lies in $\mathrm{Fl}^o$. Finally, the exponent
$e$ is \emph{twisted} on chart $k$ for the weight $\lambda$ if for every complex framed point $y$ there are
dominant weights $(a,n)$ and $(a_0,n_0)$ for which (NV) holds at $\sigma$, whose exponents differ by $e$ and
whose model weights on chart $k$ differ by $\lambda$, and a rule $g$ of weight $(a_0,n_0)$ with $g(y)\ne0$.

\begin{lemma}[Algebraicity of cuspidal sections]\label{pf:covol:lem:comb}
Let $e$ be an exponent vector with $e_0=0$ at the polarisation index, and for every chart $k$ let
$\lambda_k$ be a weight for which $e$ is twisted on chart $k$ and $P_k$ a cuspidal holomorphic section of
weight $\lambda_k$. Then there are $a_1,\dots,a_m\in\C$ and algebraic forms $f_1,\dots,f_m$ of exponent $e$
vanishing along the boundary with $P_k(x)=\sum_ia_i\,\mathrm{rd}(f_i)(k,x)$ for every chart $k$ and every
frame $x$ of it.
\formallabel{PositivityTheory.UnitarySelfGeneration.VuOnto.exists_comb}
\end{lemma}

\begin{proof}
The flag model is smooth of some relative dimension and projective over $F$
(\lref{GlobalTheory.T.toroidal_flag_smooth_projective}). For every chart $k$, $P_k$ defines a function
$h_k$ on the complex points of $\overline{\mathrm{Fl}}$ at $\sigma$, with values in the fibres of
$\mathcal L(e)$ (in the trivialisations of a covering), which reads $P_k$ at the points of the image of
chart $k$ and vanishes at all other points
(\lref{PositivityTheory.UnitarySelfGeneration.VuOnto.open_secZ}). It is holomorphic on the interior
(\lref{PositivityTheory.UnitarySelfGeneration.VuOnto.open_holoZ}, whose proof uses the twisting
hypothesis) and also at the boundary points
(\lref{PositivityTheory.UnitarySelfGeneration.VuOnto.holoOn_bdZ}, whose proof uses the cuspidality of
$P_k$). Hence $h=\sum_kh_k$ is a holomorphic
section of $\mathcal L(e)$ on $\overline{\mathrm{Fl}}_\sigma(\C)$ which vanishes outside the interior. The
images of different charts are disjoint, by the definition of a uniformisation and because the Borel
changes of frame form a groupoid (\lref{GlobalTheory.Reg.BorelGroupoid.frBor_groupoid}); so $h$ reads
$P_k$ on chart $k$. As
$\overline{\mathrm{Fl}}_\sigma$ is smooth and projective, $h$ is the section of an algebraic section of
$\mathcal L(e)$ over $\C$ (\lref{GlobalTheory.AnRead.bijective_valC}; this is GAGA). An algebraic section
over $\C$ of a line bundle on a scheme locally of finite type over $F$ which vanishes outside an open set $O$
is a finite combination $\sum_ia_is_i$ with $a_i\in\C$ and sections $s_i$ over $F$ whose non-vanishing loci
lie in $O$ (\lref{GlobalTheory.T.complex_section_finite_combination}). The sections $s_i$ (for
$O=\mathrm{Fl}^o$) define algebraic forms $f_i$ of exponent $e$ vanishing along the boundary, and reading
$h=\sum_ia_is_i$ through the uniformisation gives $P_k=\sum_ia_i\,\mathrm{rd}(f_i)(k,\cdot)$
(\lref{GlobalTheory.T.holomorphic_extension_section_reading}).
\end{proof}

\begin{proposition}[Holomorphic sections are combinations of algebraic forms]\label{pf:covol:prop:lum}
Assume (NV) and $\mathcal P$. Let $Y$ be a good deep cover as in \cref{pf:covol:def:small} with
$[E:\Q]/2>2$, $Z$ a component, $\mathrm{TF}$ a toroidal flag datum and $u$ a weight such that
$(Z,u)$ satisfies the first three conditions of admissibility in \cref{pf:covol:def:small} (only the
generation of $\Omega_Z$ is used), and let $\sigma$ be a good embedding. Then every cuspidal holomorphic
section $P$ of the model weight of $G_u$ at $\sigma$ agrees on the frames over $Z$ with $\sum_ia_i\,\mathrm{rd}(f_i)$ for some $a_i\in\C$ and algebraic forms $f_i$ of exponent
$e_u$ (the exponent of $G_u$) vanishing along the boundary of $\mathrm{TF}$.
\formallabel{PositivityTheory.Xj7N.local_unit_map_surjective}
\end{proposition}

\begin{proof}
\emph{The weight is dominant.} Let $\tau$ and $m$ with $m+1<\mathrm{mult}(\tau)$. If
$u_{\tau,m+1}>u_{\tau,m}$, every rule of exponent $e_u$ vanishes
(\lref{GlobalTheory.NonDomVanish.eq_zero_of_lt}, which states that a rule whose exponent increases
strictly at two consecutive indices of the same $\tau$ is $0$), so the evaluation map of $\mathcal C(Z,G_u)$
is $0$, which contradicts the hypothesis. Hence each $i\mapsto u_{\tau,i}$ is non-increasing.

\emph{Twisting.} Good deep covers carry a key section, a rule which does not vanish at the relevant points
(\lref{GlobalTheory.HSApp.VanishingBound.keySec}). With it, (NV) and the monotonicity just proved, the
exponent $e_u$ is twisted on every chart for the model weight of $G_u$
(\lref{PositivityTheory.UnitarySelfGeneration.VuOnto.tw_of_dec}). As $p+q\ge1$ and $e_u$ vanishes at the
polarisation index, \cref{pf:covol:lem:comb} applies to the sections $P$ on all charts and gives the claim.
\end{proof}

\begin{proposition}[Archimedean bound for small forms]\label{pf:covol:prop:archsmall}
Assume (NV) and $\mathcal P$. Let $Y$ be a good deep cover with $[E:\Q]/2>2$, let $(Z,u)$ be admissible,
$G=G_u$, and let $(b,g,A)$ be a frame of $(\nu(Z,G),\Lambda(Z,G))$. Let $\sigma\colon F\to\C$, let
$\rho\colon\Omega_Z\to\C$ be a field embedding extending $\sigma$, $\tau>0$ and $t\in\R$. Then for every
$f\in\covolSmall(Z,G,\tau,t)$
\[
|\rho(\covolGen_Zf)|\le e^{\tau t}\,\bigl|A_\sigma^{\mathsf T}(\rho(\covolGen_Zb_j))_j\bigr|^\tau .
\]
\formallabel{PositivityTheory.Xj7N.weil_height_archimedean_contribution}%
\formallabel{PositivityTheory.Xj7N.norm_small_le}
\end{proposition}

\begin{proof}
Write $S=|A_\sigma^{\mathsf T}(\rho(\covolGen_Zb_j))_j|$. Since $p+q=r\ge2$, the flag variety has positive
dimension and a toroidal flag datum exists (\lref{DifferentiableLines.nonempty_toroidal_flag_bundle}).

\emph{Step 1: spanning.} At every good embedding $\sigma'$, every holomorphic section $P$ of the model weight
of $G$ agrees on the frames over $Z$ with a complex linear combination of cusp forms of weight $G$ supported
on $Z$, read through the uniformisation: by \cref{pf:covol:prop:lum}, $P=\sum_ia_i\,\mathrm{rd}(f_i)$ there,
and for every algebraic form $f_i$ vanishing along the boundary there is a cusp form $g_i$ of weight $G$
supported on $Z$ with the same values at all frames over $Z$ at all good embeddings
(\lref{GlobalTheory.T.extVan_section_has_cusp_form_match}).

\emph{Step 2: a frame point for $\rho$.} There are a chart $k$ of the uniformisation at the good embedding
attached to $\rho$ and a frame $y_0$ of it which lies over $Z$ and is a Borel transform of the image of the
generic framed point of $Z$ under $\rho$ (\lref{GlobalTheory.CompHS.InfFr.nonempty}). Consequently there
is $c_G>0$ with $|f'(y_0)|_G=c_G\,|\rho(\covolGen_Zf')|$ for all rules $f'$ of the exponent of $G$, where
$|\cdot|_G$ is the pointwise norm of weight $G$, and $c_{\tau G}=c_G^\tau$ (the pointwise norm transforms
under the Borel change by a positive factor which is multiplicative in the weight).

\emph{Step 3: the Bergman function.} By the definition of frames (\cref{def:covolume:slopes}), the Hermitian
form of $\nu_\sigma(Z,G)$ is diagonalised by $(b,g,A)$, and Cauchy--Schwarz gives
$|\rho(\covolGen_Zf')|\le S\,\nu_\sigma(f')$ for every cusp form $f'$ of weight $G$ supported on $Z$
(\lref{PositivityTheory.Xj7N.norm_genVal_le}). Hence the Bergman function of these cusp forms at $y_0$,
$\sup_{f'\ne0}|f'(y_0)|_G^2/\nu_\sigma(f')^2$, is at most $(c_GS)^2$
(\lref{GlobalTheory.CompHS.bergCusp_le_iff}). By Step 1 it equals the Bergman function $B_G(y_0)$ of
the holomorphic sections of the model weight used in \cref{pf:covol:def:small}(c)
(\lref{GlobalTheory.CompHS.bergCusp_eq_bergOn}); so $B_G(y_0)\le(c_GS)^2$.

\emph{Step 4.} As $f$ is small, $|f(y_0)|_{\tau G}^2\le e^{2\tau t}B_G(y_0)^\tau$, that is
$c_G^{2\tau}|\rho(\covolGen_Zf)|^2\le e^{2\tau t}(c_GS)^{2\tau}$. Dividing by $c_G^{2\tau}>0$ and taking
square roots gives the claim.
\end{proof}

\begin{remark}\label{pf:covol:rem:archsmall}
Some hypotheses of the formal statements of \cref{pf:covol:prop:lum,pf:covol:prop:archsmall} are not used.
In \cref{pf:covol:prop:lum}, the Hodge package, the condition $[E:\Q]/2>2$ and the frame and spanning
conditions of admissibility are not used. In \cref{pf:covol:prop:archsmall}, of the admissibility of $(Z,u)$
only the generation of $\Omega_Z$ is used (besides the frame $(b,g,A)$, which is given separately); the
absence of base locus, the positivity of the image degree and the condition $[E:\Q]/2>2$ are not used, and
the Hodge package is used only to obtain a toroidal flag datum.
\end{remark}

\begin{proof}[Proof of \cref{pf:covol:thm:small}]
Write $G=G_u$, $\Lambda=\Lambda(Z,G)$, $\covolGen=\covolGen_Z$. Some $w_1\in\Lambda$ has
$\covolGen(w_1)\ne0$: otherwise the evaluation map vanishes on $\Lambda$, hence on $\mathcal C(Z,G)$, and its
image degree $\deg(Y,u,Z)$ would be $0$. Since the sections $s(w)$ have no common zero, there is
$w\in\Lambda$ with $\eta_X\in D(s(w))$; then $\covolGen(w)\ne0$, since otherwise the ratio condition
$\psi(\covolGen\,w_1)s(w)=\psi(\covolGen\,w)s(w_1)$ would force $\psi(\covolGen\,w_1)s(w)=0$ at $\eta_X$.
Let $R$ be the map of \cref{pf:covol:lem:ratsmall} for $s_0=s(w)$ and $\tau=8i$.

Let $f\in\covolSmall(Z,G,8i,t)$ and $h_f=\psi(\covolGen\,f)\psi(\covolGen\,w)^{-8i}$. The rational section
$h_fs(w)^{8i}$ is regular at every point of $X$ by \cref{pf:covol:lem:regular} with $\beta=\covolGen(f)$.
Its hypothesis is supplied by \lref{ChowFormHeightComparison₂.RegularPointExistence.fin_small_gen}
and \lref{ChowFormHeightComparison₂.SmallGenerators.fin_small_gen}, which state that for every
valuation $\mu$ of $\Omega_Z$ with $\mu\le1$ on $\OO_F$ and $\mu<1$ on $\mathfrak l$ there is $l\in\Lambda$
with $\mu(\covolGen f)\le\mu(\covolGen l)^{8i}$, as a consequence of the conditions at the valuations attached
to $\mathfrak l$ in the definition of small forms; the resulting regularity statements are
\lref{ChowFormHeightComparison₂.regulator_lower_bound} (points of the generic fibre) and
\lref{ChowFormHeightComparison₂.height_lower_bound_of_regulator} (all points). Next, \cref{pf:covol:prop:archsmall} with $\tau=8i$ shows that $\beta=\covolGen(f)$ satisfies
\eqref{pf:covol:eq:genbound} for every frame of $(\nu(Z,G),\Lambda)$. By \cref{pf:covol:lem:dense} with
$\beta=\covolGen(f)$ and $\tau=8i$, $\|h_fs(w)^{8i}(y)\|^2\le e^{16it}$ at all complex points. By
\cref{pf:covol:lem:ratsmall} there is a small section $S_f$ of $\overline L(t)^{\otimes8i}$ with
$R(S_f)=h_f$.

The map $f\mapsto S_f$ is injective: $R(S_f)=R(S_{f'})$ gives $\psi(\covolGen f)=\psi(\covolGen f')$ because
$\psi(\covolGen w)\ne0$, hence $\covolGen(f-f')=0$ as $\psi$ is injective, and
a cusp form of weight $8iG$ supported on $Z$ whose generic value at $Z$ vanishes is zero: its generic
values at the other components vanish because it is supported on $Z$, and \cref{pf:covol:lem:secz}
applies. This proves the inequality of cardinalities.
\end{proof}

% ---------------------------------------------------------------------------
\subsection{Non-vanishing of automorphic forms at interior points}\label{sec:pf:covol:nv}

\begin{definition}[Algebraic automorphic forms of a given weight]\label{pf:covol:def:sec}
Let $E$ be a CM field with a host datum, $\mathrm{Dw}$ a unitary datum of signature $(p,q)$ with
$p+q\ge1$, $l$ a level, $X^o$ the unitary moduli space and $c$ a coefficient datum with field $F$. For an
exponent vector $a\colon\Sigma_E\times\N\to\N$ and $n\in\N$, $\mathrm{sec}(a,n)$ is the $F$-space of
algebraic automorphic forms (rules, as in \cref{pf:covol:def:small}) of weight $(a,n)$, and $(a,n)$ is
\emph{dominant} if each $i\mapsto a(\tau,i)$ is non-increasing and $n$ is even. For an embedding
$\sigma\colon F\to\C$, a \emph{complex framed point} $y$ of $X^o$ over $\sigma$ is a complex point of $X^o$
with a frame, and $f(y)\in\C$ is the value of a rule $f$ at $y$. The \emph{non-vanishing statement} (NV)
is: for all such data, every dominant $(a,n)$ and every complex framed point $y$ there are a dominant
$(a',n')$, $g\in\mathrm{sec}(a',n')$ and $f\in\mathrm{sec}(a+a',n+n')$ with $g(y)\ne0$ and $f(y)\ne0$.
\formal{Formalism.L2.Pj.Dom}\formal{Formalism.L2.Pj.sec}\formal{Formalism.L2.Pj.evR}%
\formallabel{PositivityTheory.Xj7N.NV}
\end{definition}

\begin{theorem}[Non-vanishing at interior points]\label{pf:covol:thm:nv}
The Siegel toroidal input and the unitary toroidal input of \cref{tab:covolume:inputs} imply (NV).
\formallabel{GlobalTheory.CompanionTheory.Ar.Rec.exists_nonvanishing_of_toroidal}
\end{theorem}

\begin{proof}
Fix the data, a dominant $(a,n)$ and $y$. The unitary toroidal input, with a polarised fan provided by
the Siegel toroidal input, gives a smooth projective toroidal model $\overline Y$ of $X^o$ over $F$ with a
boundary divisor $\Delta$, given by an effective Cartier divisor whose canonical section $s_\Delta$ does not
vanish on the interior, and an integer $k\ge1$ such that $A=\omega^{\otimes k}\otimes\OO(-\Delta)$ is
ample, where $\omega$ is the determinant of the Hodge bundle of the model. Over $\overline Y$ there is an
integral flag model $\pi\colon\mathcal F\to\overline Y$ carrying automorphic line bundles $\mathcal L(e)$ for
exponent vectors $e$, with $\mathcal L(e+e')=\mathcal L(e)\otimes\mathcal L(e')$ in the Picard group: by
\lref{GlobalTheory.T.nonempty_integralFlagModel_of_smooth_coefficientField}, which states that a
smooth toroidal model on which the push-forward of the polarisation line is invertible has an integral flag
model, since the model has a theta structure by
\lref{GlobalTheory.ThetaNormal.hasTheta_of_smooth}, which states that every smooth toroidal model has
one. For the vector $w$ with all entries $1$, $\mathcal L(w)\cong\pi^*\omega$. For
$e_0(\tau,i)=\mathrm{mult}(\tau)-i$, the bundle $\mathcal L(e_0)$ is ample on $\pi^{-1}(V)$ for the members
$V$ of an affine open cover of $\overline Y$, by
\lref{GlobalTheory.T.borelWeil_flagModel_detInverse_isAmple_local}, which states this ampleness
(the Borel--Weil property of the flag model). Hence $\mathcal L(e_0)$ is ample relative to $\pi$, and
$P=\pi^*A^{\otimes m}\otimes\mathcal L(e_0)$ is ample for all large $m$ \cite[\S1.7]{Lazarsfeld2004};
fix such $m\ge1$. Put $\tau_0=mk\,w+e_0$; then $\mathcal L(\tau_0)=P\otimes\pi^*\OO(\Delta)^{\otimes m}$ in the
Picard group, and the section $\pi^*s_\Delta^m$ of the second factor does not vanish at the point $z$ of
$\mathcal F$ defined by the framed point $y$, which lies over the interior.

For every line bundle $\mathcal N$ on $\mathcal F$, the sheaf $\mathcal N\otimes P^{\otimes u}$ is generated by
global sections for all large $u$ \cite[Thm.~1.2.6]{Lazarsfeld2004}; multiplying by $\pi^*s_\Delta^{mu}$
gives, for all large $u$, sections of $\mathcal N\otimes\mathcal L(u\tau_0)$ not vanishing at $z$. Apply this
to $\mathcal N=\mathcal L(0)$ and to $\mathcal N=\mathcal L(e_A)$, where $e_A$ is the exponent vector of the
weight $(a,n)$, and take a common $u\ge1$: we obtain sections $s_1$ of $\mathcal L(u\tau_0)$ and $s_2$ of
$\mathcal L(e_A+u\tau_0)$ not vanishing at $z$. Let $a'(\tau,i)=u\,(mk+\mathrm{mult}(\tau)-i)$ and $n'=0$;
then $(a',0)$ is dominant, its exponent vector is $u\tau_0$, and the exponent vector of $(a+a',n)$ is
$e_A+u\tau_0$. The global sections $s_1,s_2$ define rules $g\in\mathrm{sec}(a',0)$ and
$f\in\mathrm{sec}(a+a',n)$, whose values at $y$ are the values of $s_1,s_2$ at $z$ in the frame given by $y$,
hence nonzero.
\end{proof}

% ---------------------------------------------------------------------------
\subsection{Statements from the literature used by the assembly}\label{sec:pf:covol:literature}

The following statements, proved in the development, enter the proof of \cref{thm:covol} directly. They
are published results (or immediate consequences of published inputs), and we only recall them.

\begin{littheorem}[Green {\cite{Green1955}}]\label{pf:covol:lit:green}
The parametrisation of the irreducible cuspidal characters of $\GL_d(\F_g)$ by regular characters of
$\F_{g^d}^\times$ holds (\cref{thm:literature:green}); consequently $\sfname{Green}$
(\cref{def:green}) holds, by \cref{lit:top:green}.
\formallabel{GlobalTheory.greenRegularCharacters_holds}
\end{littheorem}

\begin{littheorem}[Fubini--Study heights of subvarieties of projective space {\cite{Philippon1991,Soule1991}}]\label{pf:covol:lit:fs}
Assume $\mathcal P$. Let $Y\subseteq\mathbb P^N_{\OO_K}$ be an integral closed subscheme which is an
arithmetic variety of relative dimension $e$, let $A_\sigma\in\GL_{N+1}(\C)$ with
$A_{\bar\sigma}=\overline{A_\sigma}U_\sigma$, $U_\sigma$ unitary, and let $f$ be an integral Chow form of
dimension $e$ and degree $\delta$ of the generic ideal of $Y$. Then the restriction to $Y$ of $\OO(1)$
with the Fubini--Study metrics of the coordinate sections and the matrices $A_\sigma$ has
self-intersection $[K:\Q]\,\covolHform(\delta,1,A,f)$. (The development derives this from the arithmetic
intersection theory contained in $\mathcal P$.)
\formallabel{FubiniStudyChowHeightFormula₂.fubiniStudy_chowHeight_formula}
\end{littheorem}

\begin{littheorem}[Integral Chow forms {\cite{ChowVanDerWaerden1937}}]\label{pf:covol:lit:chowint}
Let $k$ be a field of characteristic $0$ and $P\subseteq k[X_1,\dots,X_N]$ a homogeneous prime with
$\dim k[X]/P=e+1$. Then $P$ has an integral Chow form of dimension $e$ whose degree is the degree of $P$
(the normalised leading coefficient of its Hilbert polynomial), and every integral Chow form of $P$ of
dimension $e$ and the same degree is a nonzero multiple of it.
\formallabel{GlobalTheory.Published.chowFormIntExists_holds}
\end{littheorem}

\begin{littheorem}[Chow forms of hyperplane sections {\cite[\S1]{Philippon1986}}]\label{pf:covol:lit:cut}
Let $k$ be a field of characteristic $0$, $P\subseteq k[X_1,\dots,X_N]$ a homogeneous prime with
$\dim k[X]/P=e+2$, $u\in k^N$ with $\ell_u\notin P$, and $f$ an integral Chow form of $P$ of dimension
$e+1$. If $g_Q$ are integral Chow forms of dimension $e$ of the components $Q$ of the cycle
$[P+(\ell_u)]$, then $f(u,\cdot)=c\prod_Qg_Q^{m_Q}$ with $c\in k^\times$ and $m_Q$ the multiplicities.
(The development deduces this from \cref{pf:covol:lem:chowcalc}(a) and the uniqueness of Chow forms up
to scalars.)
\formallabel{ChowFormCompatibility.chow_form_hyperplane_section}
\end{littheorem}

\begin{littheorem}[Toroidal charts {\cite{FaltingsChai1990,Lan2013}}]\label{pf:covol:lit:tor}
(a) The unitary toroidal input of \cref{tab:covolume:inputs} implies the existence of \emph{toroidal charts
for a given fan}: for every admissible polarised fan without self-intersection, the unitary moduli spaces
have smooth toroidal models with simple normal crossings boundary, charts at the boundary points and an
ample polarisation (the Kodaira--Spencer clause of the input is dropped). (b) This statement and the
Siegel toroidal input imply the existence of \emph{toroidal charts}: smooth toroidal models with simple
normal crossings boundary and charts at all boundary points.
\formallabel{PositivityTheory.Prw.toroidalChartsForFan_of}\formallabel{GlobalTheory.Published.toroidalCharts_of_siegel_of_forFan}
\end{littheorem}

In (b), the polarised fan of the Siegel toroidal input in genus $[E:\Q]q+1$ is refined to a fan which is
invariant under the CM structure, polarised and without self-intersection, by
\lref{GlobalTheory.PolFan.exists_fine_invariant_fan}, which states that every polarised fan for
$\Z$-lattices of rank $\ge[E:\Q]q$ restricts and refines to such a fan for projective $\OO_E$-modules of
rank $\le q$; this fan is then used in (a).

\begin{littheorem}[Lengths of proper intersections {\cite[\S7.1]{Fulton1998}}]\label{pf:covol:lit:length}
Let $X$ be an integral scheme, projective of dimension $D$ over a field $k$, $L$ a line bundle on $X$ and
$s_1,\dots,s_D$ sections of $L$ whose common zero scheme $C$ is finite with Cohen--Macaulay local rings.
Then the sections form a regular sequence at the points of $C$, $C$ is finite over $k$, and the top
self-intersection of $L$ equals $\sum_{z\in C}[\kappa(z):k]\,\len(\OO_{C,z})=\dim_k\Gamma(C,\OO_C)$.
(In the development this statement, which belongs to the package $\mathcal Q$ of
\cref{pf:covol:def:packages}, is supplied by a theorem whose hypothesis, the Hodge package, is not used.)
\formal{GlobalTheory.CovTop.PublishedHodgeBundleTheorems.q12}
\end{littheorem}

% ---------------------------------------------------------------------------
\subsection{The final assembly}\label{sec:pf:covol:assembly}

We use the families, points and constants of \cref{sec:covolume:statement} and the lift predicates
$\mathsf{GramAt}$, $\mathsf{RecAt}$, $\mathsf{CutAt}$, $\mathsf{BadAt}$, $\mathsf{LamAt}(A)$,
$\mathsf{SizeAt}(PP,A)$, $\mathsf{CoreAt}(C)$, $\mathsf{LowAt}(C)$ of \cref{def:covolume:predicates}. For a
point $p$ at $(x,\iota)$ we write $b$, $e_z$ for the parameters of its size data,
$\delta=\exp(-x^{e_z/7})$, $\covolMainT_i=\covolMainT_i(p)$ and $U=\covolUnitE(p)$. Recall that the lift
condition attached to a predicate $Q(p,i,PP,C)$ is
\begin{multline*}
\mathrm{Closed}\Bigl(\mathrm{FamHyp},(F,C_H)\mapsto\exists C\ge0:\\
\mathrm{EventuallyAtAllPoints}\bigl(F,p\mapsto\forall i\in J:\ Q(p,i,PP(x,\iota),C)\bigr)\Bigr),
\end{multline*}
where $PP(x,\iota)$ is a parameter of the family at $(x,\iota)$; the eight lift conditions are those
attached to the eight predicates above, the existential constant playing the role of $A$ for
$\mathsf{LamAt}$ and of the second argument for $\mathsf{SizeAt}$, $\mathsf{CoreAt}$ and $\mathsf{LowAt}$.

\begin{lemma}[Covolume bound at a point]\label{pf:covol:lem:lcvle}
Let $p$ be a point at $(x,\iota)$ with $x>0$, let $i\in J$, $A,B,PP\in\R$ and $C_1,C_2\ge0$ with
\begin{equation}\label{pf:covol:eq:loss}
12\,e^{-x^{e_z/7}}\bigl(x^A+x^B+C_1+C_2+1\bigr)\le1 .
\end{equation}
If $\mathsf{GramAt}$, $\mathsf{RecAt}$, $\mathsf{CutAt}$, $\mathsf{BadAt}$, $\mathsf{LamAt}(A)$,
$\mathsf{SizeAt}(PP,B)$, $\mathsf{CoreAt}(C_1)$ and $\mathsf{LowAt}(C_2)$ hold at $(p,i)$, then
\[
\frac1D\,\covolLcv(X_{W,i},M_i)\le\covolMainT_i+(C_1+C_2+1)\,U .
\]
\formal{GlobalTheory.CovTop.Assembly.lcv_le}
\end{lemma}

\begin{proof}
Let $Q=(x^A+x^B+C_1+C_2+1)U$, let $\Lambda^{\mathrm{type}}$, $\Lambda^{\mathrm{all}}$,
$\Lambda^{\mathrm{good}}$ be the lattices of \cref{sec:covolume:assembly}, $e=e_i$ and $g=|G_p|$. By
$\mathsf{SizeAt}(PP,B)$, $1\le b\le U$ and $|\covolMainT_i|\le x^Bb$; in particular $U\ge0$ and $Q\ge U$.
We use the estimate for the slope of the good part (\cref{lem:covolume:slope-good}, formally
\lref{GlobalTheory.CutNum.slope_good_ge}): if the norms have a Gram frame, $L\subseteq M$ with $M$
finitely generated and $L\ne0$, $e$ is a Pythagorean idempotent, $g\ge1$ is an integer with
$g\,e(L)\subseteq L$, $0\le\delta\le\frac1{12}$, $\rank(L\cap\ker e)\le3\delta\rank L$,
$\mu_{\max}(M)\le Q$, $-Q\le\mu(L)$ and $2\log g\le Q$, then $\mu(L\cap\operatorname{im}e)\ge\mu(L)-12\delta Q$.
Apply it to $L=\Lambda^{\mathrm{type}}\subseteq M=\Lambda^{\mathrm{all}}$. The hypotheses hold: the Gram
frame, the finite generation of $M$ and $L\ne0$ (because $0\ne\Lambda^{\mathrm{good}}\subseteq L$) come
from $\mathsf{GramAt}$; $e^2=e$, $e$ Pythagorean and $g\,e(L)\subseteq L$ from $\mathsf{CutAt}$;
$12\delta\le1$ from \eqref{pf:covol:eq:loss}, because $x^A,x^B>0$ and $C_1+C_2+1\ge1$; the rank condition
is $\mathsf{BadAt}$; $\mu_{\max}(M)\le x^Ab\le x^AU\le Q$ by $\mathsf{LamAt}(A)$; by $\mathsf{CoreAt}(C_1)$,
$\mathsf{LowAt}(C_2)$ and $\mathsf{SizeAt}$,
\[
\mu(L)\ge\mu_{\mathrm{core}}-C_1U\ge-\covolMainT_i-(C_1+C_2)U\ge-x^BU-(C_1+C_2)U\ge-Q;
\]
and $2\log g\le b\le U\le Q$ by $\mathsf{CutAt}$. Hence $\mu(\Lambda^{\mathrm{good}})\ge\mu(L)-12\delta Q$.
Multiplying \eqref{pf:covol:eq:loss} by $U\ge0$ gives $12\delta Q\le U$. With $\mathsf{RecAt}$,
$\mathsf{CoreAt}(C_1)$ and $\mathsf{LowAt}(C_2)$,
\[
\frac1D\covolLcv(X_{W,i},M_i)\le-\mu(\Lambda^{\mathrm{good}})\le-\mu(L)+U\le-\mu_{\mathrm{core}}+(C_1+1)U
\le\covolMainT_i+(C_1+C_2+1)U.\qedhere
\]
\end{proof}

\begin{lemma}\label{pf:covol:lem:loss}
For $c>0$ and $A,B,C\in\R$ there is $x_0$ such that $12\,e^{-x^c}(x^A+x^B+C)\le1$ for all $x\ge x_0$.
\formal{GlobalTheory.CovTop.Assembly.exists_loss_le}\formal{GlobalTheory.CovTop.Assembly.tendsto_exp_mul_rpow}
\end{lemma}

\begin{proof}
For every $a\in\R$, $e^{-x^c}x^a=e^{-y}y^{a/c}\to0$ as $x\to\infty$, where $y=x^c\to\infty$. Hence
$12e^{-x^c}(x^A+x^B+Cx^0)\to0<1$.
\end{proof}

\begin{theorem}[\cref{thm:covolume:assembly}]\label{pf:covol:thm:assembly}
The eight lift conditions of \cref{def:covolume:predicates} imply $\sfname{CovolumeBoundOfX}$.
\formal{GlobalTheory.CovTop.Assembly.covolX_of_lift_bounds_and_coreLe}
\end{theorem}

\begin{proof}
A finite conjunction of statements of the form $\mathrm{Closed}(\mathrm{FamHyp},P_k)$ is again of this form
for the conjunction of the $P_k$: take the least of the thresholds $\eps_1$ and a transported constant
dominating those of the $P_k$; and if $P$ implies $P'$ for all $\eps$, families $F$ and constants $C_H$
satisfying the conditions of $\mathrm{Closed}$, then $\mathrm{Closed}(\mathrm{FamHyp},P)$ implies
$\mathrm{Closed}(\mathrm{FamHyp},P')$. Similarly a finite conjunction of statements
$\mathrm{EventuallyAtAllPoints}(F,Q_k)$ holds with the largest of the thresholds $x_0$.

So fix $\eps$, a family $F$ and $C_H$ as in $\mathrm{Closed}$, and the constants provided by the eight
conditions: $A$ for $\mathsf{LamAt}$, $B$ for $\mathsf{SizeAt}$, $C_1\ge0$ for $\mathsf{CoreAt}$ and
$C_2\ge0$ for $\mathsf{LowAt}$; let $x_1'$ be a common threshold for the eight
$\mathrm{EventuallyAtAllPoints}$ statements. The exponent $e_z$ of the size data is the same at all
$(x,\iota)$ and is positive, since it is positive at the first index (the scale axioms); by
\cref{pf:covol:lem:loss} with $c=e_z/7$ and $C=C_1+C_2+1$ there is $x_0'$ such that
\eqref{pf:covol:eq:loss} holds for $x\ge x_0'$. Put $C_0=C_1+C_2+1\ge0$ and take the threshold
$\max(x_0',x_1')$. Let $x$ exceed it (so that also $x>x_1>1$, where $x_1$ is the threshold of the
family), $\iota\in I_x$, $p$ a point at $(x,\iota)$ satisfying $\mathrm{PointConditions}'$ and
$\mathrm{RecTied}$, $C\ge C_0$ and $i\in J$. The eight predicates hold at $(p,i)$, so
\cref{pf:covol:lem:lcvle} gives $\frac1D\covolLcv(X_{W,i},M_i)\le\covolMainT_i+C_0U$. Since $U\ge b\ge1>0$ by
$\mathsf{SizeAt}$, this is at most $\covolMainT_i+CU$, which is $l_3(i)$ for the numerical datum
$p.N(C,C_H)$; this is the inequality required by $\sfname{CovolumeBoundOfX}$ (\cref{def:covolX}).
\end{proof}

\begin{remark}\label{pf:covol:rem:assembly}
The formal statements of this subsection belong to the formal proof of the lift estimates
(\cref{thm:covolume:lift}), whose last step is \cref{pf:covol:thm:assembly}, applied through a
restatement of it. The constant $C_H$ enters only through the constants and thresholds provided by the
lift conditions.
\end{remark}

% ---------------------------------------------------------------------------
\subsection{Proof of the covolume bound}\label{sec:pf:covol:main}

We first record the two remaining notions which occur in the statements used.

\begin{definition}[The four packages]\label{pf:covol:def:packages}
\begin{enumerate}[label=(\alph*)]
\item The \emph{unitary geometric package} $\mathcal H_X$ consists of the unitary toroidal input and the
  Lie seminegativity and point separation inputs of \cref{tab:covolume:inputs}, of $\sfname{Green}$
  (\cref{def:green}), and of the formula for Fubini--Study heights of subvarieties of projective space
  (\cref{pf:covol:lit:fs}).
\item The \emph{toroidal package} $\mathcal H_H$ consists of the Siegel toroidal input, the existence of
  toroidal charts for a given fan and of toroidal charts (\cref{pf:covol:lit:tor}), and $\sfname{ArDegree}$.
\item The \emph{volume package} $\mathcal V$ consists of the \emph{self-intersection hypothesis}: $\mathcal P$
  implies the conclusion of \cref{pf:covol:thm:comparison} for all data; and the \emph{small-sections
  hypothesis}: $\mathcal P$ implies the conclusion of \cref{pf:covol:thm:small} for all data.
\item The \emph{quoted package} $\mathcal Q$ consists of the Kawamata--Viehweg vanishing theorem, the
  $L^2$ extension theorem for log canonical pairs, GAGA for holomorphic sections, the numerics of nef
  divisors and the existence of Chow forms (all statements of $\mathcal P$), the existence of integral
  Chow forms (\cref{pf:covol:lit:chowint}), Chow forms of hyperplane sections (\cref{pf:covol:lit:cut}) and
  the length formula for proper intersections (\cref{pf:covol:lit:length}).
\end{enumerate}
\formallabel{GlobalTheory.CovTop.HypothesesX}\formallabel{GlobalTheory.CovHt.HypothesesH}%
\formal{SelfIntersectionVolumeEstimates.VolumeEstimateHypotheses}\formal{GlobalTheory.HSCore.Quoted}
\end{definition}

\begin{definition}[Mean slopes on deep covers]\label{pf:covol:def:deepstr}
Let $Y$ be a deep cover as in \cref{pf:covol:def:small}, of a datum of signature $(p,q)$ with CM type
$\Psi$, and let $u=(u_{\tau,a})$ be a weight of the datum.
\begin{enumerate}[label=(\alph*)]
\item For $g>0$, $u$ has \emph{gap} $g$ if every $a\mapsto u_{\tau,a}$ is strictly decreasing and
  $|u_{\tau,a}-u_{\tau,a'}|\ge g$ for $a\ne a'$ with $a\ge p$. Put
  $\mathrm{bnd}(u,g)=4|\Psi|(p+q)^2\bigl(1+\max_{\tau,a}|u_{\tau,a}|/g\bigr)$ and, for $s\in\R$,
  $\mathrm{err}_s(u,g)=(\log\mathrm{bnd}(u,g)+s)\|u\|_1+q\sum_\tau(\max_au_{\tau,a}-\min_au_{\tau,a})$.
\item For $\delta>0$ and a function $\mathrm{vol}\colon\R\to[0,\infty]$, the \emph{volume slope} is
  $\mathrm{volSlope}(\delta,\mathrm{vol})=\inf_{t\in\R}\bigl(\mathrm{vol}(t)/\delta-t\bigr)\in[-\infty,\infty]$.
  For $D,d,a\in\N$ and a sequence $m(\tau)\in\N\cup\{\infty\}$, the \emph{counting volume} is
  $\limsup_{k\to\infty}D!\,\log^+m(ak)/(d\,(ak)^{D+1})$.
\item For a family $\mathcal V$ of valuations, a component $Z$ and $a\in\N$, let $\mathrm{vol}_{a}(Z,u)(t)$
  be the counting volume, with $D=\dim\mathrm{Fl}$ and $d=[F:\Q]$, of the numbers of cusp forms of weight
  $\tau G_u$ on $Z$ which are integral for $\mathcal V$ and have pointwise norm at most $e^{\tau t}$ at all
  frame points over $Z$ at the good embeddings; and let
  $\mathrm{slope}_a(Z,u)=\mathrm{volSlope}(\gamma_ZV(u),\mathrm{vol}_a(Z,u))$, where
  $V(u)=\prod_\tau\prod_{i<l}(u_{\tau,i}-u_{\tau,l})$ is the Vandermonde product of $u$ and $\gamma_Z$ is
  the growth constant of the dimensions of the spaces of cusp forms on $Z$ for the standard weight $u_0$,
  divided by $V(u_0)$.
\item $\mathrm{MeanAll}(s,C)$ holds on $Y$ if for every weight $u$ with gap $g>0$ there are real numbers
  $\mathrm{sl}(Z)\le\mathrm{slope}_8(Z,u)$, computed with the family of valuations attached to the primes
  up to $\mathrm{bnd}(u,g)$, whose average over the components $Z$ is at least
  $-C\,\mathrm{err}_s(u,g)$.
\item $\sfname{DeepStR}$ is the statement: there is $C\ge0$ such that for every global datum $\mathcal D$
  and every core datum $T$ of $\mathcal D$ satisfying the structural conditions $\mathsf{LawsJ}$, with
  $m^2\le k$, every $1\le j<m$ and every good deep cover $Y$ of $T$ at level $j$ whose level modulus has at
  most $15$ prime factors, $\mathrm{MeanAll}(\ell(\mathcal D),C)$ holds on $Y$, where
  $\ell(\mathcal D)=\log\mathrm{rd}_E+|\frac{L'}{L}(1,\chi_d)|$.
\end{enumerate}
This makes \cref{def:covolume:deepstr} of Part I explicit.
\formallabel{GlobalTheory.CompanionTheory.Ar.Plug.DeepStR}\formallabel{GlobalTheory.CovHt.MeanAll}%
\formallabel{GlobalTheory.CovHt.MeanGe}\formallabel{GlobalTheory.CovHt.WtGap}\formallabel{GlobalTheory.CovHt.bndB}%
\formallabel{GlobalTheory.CovHt.errX}\formallabel{GlobalTheory.CovHt.ell}\formallabel{GlobalTheory.CompHt.volSlope}%
\formallabel{GlobalTheory.CompHt.countVol}\formallabel{GlobalTheory.CompHt.compVol}%
\formallabel{GlobalTheory.CompHt.compVolSlope}\formallabel{GlobalTheory.CompHt.topDeg}
\end{definition}

\begin{lemma}\label{pf:covol:lem:volslope}
For $\delta>0$, a function $\mathrm{vol}\colon\R\to[0,\infty]$ and $a\in\R$, one has
$a\le\mathrm{volSlope}(\delta,\mathrm{vol})$ if and only if $\delta(a+t)\le\mathrm{vol}(t)$ for every
$t\in\R$ with $a+t\ge0$.
\formallabel{GlobalTheory.CompHt.le_volSlope_iff}
\end{lemma}

\begin{proof}
$a\le\mathrm{vol}(t)/\delta-t$ is equivalent to $\delta(a+t)\le\mathrm{vol}(t)$, and it holds trivially when
$a+t<0$, as $\mathrm{vol}(t)\ge0$. (Formally the right-hand side is written with the truncation
$\max(\delta(a+t),0)$.)
\end{proof}

\begin{theorem}[\cref{thm:covol}]\label{pf:covol:thm:covol}
The fifteen statements of \cref{tab:covolume:inputs} imply $\sfname{CovolumeBoundOfX}$.
\formallabel{SelfIntersectionVolumeEstimates.topCovolSlim2_use_alt_alt}
\end{theorem}

\begin{proof}
Assume the fifteen statements. We verify the four packages of \cref{pf:covol:def:packages} and
$\sfname{DeepStR}$, and apply the lift estimates.

(1) $\mathcal H_X$: the unitary toroidal, Lie seminegativity and point separation inputs are among the
fifteen statements; $\sfname{Green}$ holds by \cref{pf:covol:lit:green}; and the Fubini--Study height
formula follows from $\mathcal P$ by \cref{pf:covol:lit:fs}.

(2) $\mathcal H_H$: the Siegel toroidal input and $\sfname{ArDegree}$ are among the fifteen statements,
and the two statements on toroidal charts follow from the Siegel and unitary toroidal inputs by
\cref{pf:covol:lit:tor}.

(3) $\mathcal V$: the self-intersection hypothesis is \cref{pf:covol:thm:comparison}. The small-sections
hypothesis is \cref{pf:covol:thm:small}, once (NV) is known; and (NV) holds by \cref{pf:covol:thm:nv},
applied to the Siegel and unitary toroidal inputs.

(4) $\mathcal Q$: the first five statements are part of $\mathcal P$, and the last three hold by
\cref{pf:covol:lit:chowint,pf:covol:lit:cut,pf:covol:lit:length}.

(5) By \cref{thm:deepstr}, i.e.\ \lref{PositivityTheory.HtPt.Ah.deep_slim_use}, which states that the
resolution, principalisation, Poincar\'e--Lelong, $\sfname{HSUse}$, Moriwaki, Hodge--Tate, Zarhin and nef
Riemann--Roch inputs imply $\sfname{DeepStR}$ (\cref{pf:covol:def:deepstr}), the first eight inputs give
$\sfname{DeepStR}$.

(6) By \cref{thm:covolume:lift}, i.e.\ \lref{SelfIntersectionVolumeEstimates.topCovolG_alt_alt},
which states that $\mathcal V$, $\mathcal H_X$, $\mathcal P$, $\mathcal H_H$, the Grothendieck--Messing input,
$\mathcal Q$ and $\sfname{DeepStR}$ imply $\sfname{CovolumeBoundOfX}$, the packages (1)--(4), the inputs
$\mathcal P$ and Grothendieck--Messing, and (5) give $\sfname{CovolumeBoundOfX}$.
\end{proof}

\begin{remark}\label{pf:covol:rem:inputs}
(1) The resolution, principalisation, Poincar\'e--Lelong, $\sfname{HSUse}$, Moriwaki, Hodge--Tate, Zarhin
and nef Riemann--Roch inputs are used only through $\sfname{DeepStR}$; the other seven inputs enter the
four packages and the lift estimates. (2) All fifteen inputs are theorems of the development:
$\sfname{HSUse}$ and $\sfname{MoriwakiBigStatement}$ by \cref{thm:covolume:hs}, the resolution inputs by
\cref{thm:covolume:resolution}, and the others as described in \cref{sec:literature}; hence
$\sfname{CovolumeBoundOfX}$ holds. (3) Inside the lift estimates, $\sfname{CovolumeBoundOfX}$ is obtained
from the eight lift conditions by \cref{pf:covol:thm:assembly}; the results of
\cref{sec:pf:covol:chowht,sec:pf:covol:comparison,sec:pf:covol:small,sec:pf:covol:nv} are used there to
bound numbers of small automorphic forms on deep covers by arithmetic volumes of models of their images.
\end{remark}

% ---------------------------------------------------------------------------
\subsection{Further results used in other sections}\label{sec:pf:covol:further}

This subsection proves statements of the development which are used in other sections of Part~II: an
exponent and a lemma on levels used for deep covers, GAGA for sections of line bundles, the identification
of the Hodge determinant on flag covers, the two flat line statements and the flat degree formula.

\subsubsection*{Deep covers}

\begin{definition}[Exponents of primes]\label{pf:covol:def:htex}
For the host field $E$ of a global datum and $p\in\N$ let
$\mathrm{ex}(p)=\frac1{p-1}+40\,\frac{v_p(|d_E|)}{[E:\Q]}$, where $d_E$ is the discriminant of $E$ and
$v_p$ the $p$-adic valuation (for $p=1$ the first term is read as $0$). It is used in the height comparisons
for deep covers.
\formallabel{GlobalTheory.CovHt.htEx}
\end{definition}

Recall that a level $l$ of a host datum has a modulus $M(l)$, an odd integer $\ge3$, and that $l$ has
\emph{good reduction at a prime $p$} if the prime-to-$p$ part $M(l)/p^{v_p(M(l))}$ is at least $3$.

\begin{lemma}[Good reduction and multiples of the level]\label{pf:covol:lem:goodmul}
\begin{enumerate}[label=(\alph*)]
\item A level $l$ has good reduction at a prime $p$ if and only if $M(l)$ is not a power of $p$.
\item Let $l,l_2$ be levels with $M(l_2)=b\,M(l)$ for some $b\in\N$. If $l$ has good reduction at every
  prime, then so has $l_2$.
\end{enumerate}
\formallabel{GlobalTheory.HostProbe.goodAt_iff}\formallabel{GlobalTheory.LawsDeep.goodAt_of_mul}
\end{lemma}

\begin{proof}
(a) The prime-to-$p$ part of $M(l)$ is odd (as $M(l)$ is odd), so it is at least $3$ unless it is $1$,
that is, unless $M(l)=p^{v_p(M(l))}$; and if $M(l)=p^k$ for some $k$, then the prime-to-$p$ part is $1$.
(b) By (a) it suffices to show that $M(l_2)$ is not a power of $p$. If $M(l_2)=p^k$, then $M(l)$ divides $p^k$, so $M(l)=p^i$ for some $i\le k$, and $l$ would not
have good reduction at $p$.
\end{proof}

\subsubsection*{GAGA for sections}

\begin{theorem}[GAGA for sections]\label{pf:covol:thm:gaga}
Let $X$ be a scheme, proper and smooth over $\C$, and $L$ a line bundle on $X$, given by transition
functions $g_{ij}\in\Gamma(U_i\cap U_j,\OO_X)^\times$ on an open cover $(U_i)$. Let $h_i\colon U_i(\C)\to\C$ be
functions such that
\begin{enumerate}[label=(\roman*)]
\item for every open set $B\subseteq\C^m$ and every holomorphic family $b\mapsto y(b)$ of complex points of
  $U_i$ parametrised by $B$, the function $b\mapsto h_i(y(b))$ is holomorphic on $B$; and
\item $h_i(x)=g_{ij}(x)h_j(x)$ for every complex point $x$ of $U_i\cap U_j$.
\end{enumerate}
Then there is a global section of $L$ whose local function on $U_i$ takes the value $h_i(x)$ at every complex
point $x$ of $U_i$, for every $i$.
\formallabel{GlobalTheory.CompanionTheory.Gaga.gagaSections}
\end{theorem}

The theorem is a form of Serre's GAGA \cite{Serre1956}; the development proves it by the following direct
argument.

\begin{lemma}[The equidimensional case]\label{pf:covol:lem:gagapure}
\Cref{pf:covol:thm:gaga} holds if $X$ is smooth of pure relative dimension $D$ over $\C$.
\formallabel{GlobalTheory.CompanionTheory.Gaga.gagaSections_pure}
\end{lemma}

\begin{proof}
Every complex point of $X$ lies in an \emph{algebraic chart}: a neighbourhood biholomorphic to an open set
of $\C^D$ through $D$ algebraic functions. By (i) applied to the inverses of the charts, the $h_i$ are
holomorphic in the charts, and by (ii) they form a holomorphic section $s=(s_i)$ of $L$ on $X(\C)$
(\lref{GlobalTheory.T.holomorphic_gluing_local_functions}). We show that $s$ is algebraic
(\lref{GlobalTheory.CompanionTheory.Gaga.holomorphicSections_algebraic_of_isProper}).

\emph{Local algebraicity.} Let $x\in U_i$. As $X$ is smooth, $x$ has an open neighbourhood $U$ such that
$\Gamma(V,\OO_X)$ is a domain for every open $V\subseteq U$ containing $x$
(\lref{GlobalTheory.exists_open_isDomain_of_smooth_complex}); let $V\ni x$ be an affine open subset of
$U\cap U_i$. Because $X$ is proper, $s_i$ has polynomial growth on $V(\C)$: there are regular functions
$u_1,\dots,u_n$ on $V$ and constants $C,m$ with $|s_i(y)|\le C(1+\sum_l|u_l(y)|)^m$ for $y\in V(\C)$
(\lref{GlobalTheory.T.holomorphic_section_moderate_growth}). A holomorphic function of polynomial growth
on $V(\C)$, $V$ smooth affine with a domain as ring of functions, is a quotient of regular functions: there
are $a,b\in\Gamma(V,\OO_X)$, $b\ne0$, with $b\,s_i=a$ on $V(\C)$
(\lref{GlobalTheory.SerreSections.exists_mul_eq_of_polynomial_growth}). A continuous function on $V(\C)$
which equals $a/b$ is regular: there is $c\in\Gamma(V,\OO_X)$ with $c=s_i$ on $V(\C)$
(\lref{GlobalTheory.exists_ev_eq_of_continuousOn_quotient}).

\emph{Gluing.} Regular functions on an open subset of $X$ which agree at all complex points are equal, $X$
being reduced and of finite type over $\C$
(\lref{GlobalTheory.eq_of_ev_eq_of_smoothOfRelativeDimension}). Hence the local functions $c$ agree on
overlaps and glue to $f_i\in\Gamma(U_i,\OO_X)$ with $f_i=s_i$ on $U_i(\C)$; for the same reason
$f_i=g_{ij}f_j$ on $U_i\cap U_j$. So $(f_i)$ is a global section of $L$ with the required values.
\end{proof}

\begin{proof}[Proof of \cref{pf:covol:thm:gaga}]
A scheme smooth over $\C$ is the disjoint union of open and closed subschemes $W_D$, $D\in\N$, with $W_D$
smooth of pure relative dimension $D$ (\lref{GlobalTheory.T.smooth_scheme_decompose_relative_dimension}).
Each $W_D$ is proper over $\C$, being closed in $X$, and the restrictions of $L$ and of the $h_i$ to $W_D$
satisfy (i) and (ii), since holomorphic families of points of $W_D$ are holomorphic families of points of
$X$. By \cref{pf:covol:lem:gagapure} there are global sections $\sigma_D$ of $L|_{W_D}$ with the required
values. Since the $W_D$ form a disjoint open cover, for every $i$ the local functions of the $\sigma_D$ on
the sets $U_i\cap W_D$ glue uniquely to a function $b_i$ on $U_i$
(\lref{GlobalTheory.CompanionTheory.Gaga.glue}); the relations $b_i=g_{ij}b_j$ hold because they hold on
every $W_D$, and every complex point of $X$ lies in some $W_D$, which gives the values.
\end{proof}

\subsubsection*{The Hodge determinant on flag covers}
Let $Y$ be a model of $X^o$ over $F$ and $\pi\colon\mathrm{Fl}\to Y$ a \emph{flag cover}: $\mathrm{Fl}$ is a
normal scheme, proper over $Y$, with a \emph{frame system}, which consists of an open cover $(U_k)_k$ of
$\mathrm{Fl}$, charts $M_k$ of $Y$ (each providing a basis of the Lie algebra of the semi-abelian scheme of
$Y$ over an open set containing $\pi(U_k)$, of size $n$; on overlaps the two bases differ by a matrix
$\mathrm{chg}_{kk'}$), and $n\times\mathrm{mult}(\tau)$ matrices $u_k(\tau)$ of functions on $U_k$, with a
left inverse, such that on $U_k\cap U_{k'}$
\[
\mathrm{chg}_{kk'}\,u_{k'}(\tau)=u_k(\tau)\,B_{kk'}(\tau)
\]
with upper triangular matrices $B_{kk'}(\tau)$. The columns of $u_k(\tau)$, read in the basis of chart $k$,
are the \emph{flag vectors} $v_k(\tau,i)$, $i<\mathrm{mult}(\tau)$. Line bundles on $\mathrm{Fl}$ are given by
cocycles on $(U_k)$ and compared through their classes. The \emph{Hodge determinant}
$\det\Omega$ is given by the cocycle $(\det\mathrm{chg}_{kk'})^{-1}$; for an exponent vector $e$, the
automorphic line bundle $\mathcal L(e)$ is given by the cocycle $\prod_\iota T_{kk'}(\iota)^{-e(\iota)}$,
where $T_{kk'}(\tau,i)$ is the diagonal entry $B_{kk'}(\tau)_{ii}$ and $T_{kk'}$ at the polarisation index is
a unit attached to the two charts (it plays no role below); and $\mathcal L_\lambda=\mathcal L(e_\lambda)$ for a weight
$\lambda$ with exponent vector $e_\lambda$. For the weight $\mu_0$ (equal to $1$ on the first $p$ entries
and to $-1$ on the others), $e_{\mu_0}$ is $0$ at the polarisation index and $1$ at all $(\tau,i)$
(\lref{GlobalTheory.FlagMod.unitaryWeightExponentVector_signatureWeightMuZero}).

\begin{proposition}[The Hodge determinant on flag covers]\label{pf:covol:prop:lbiso}
On every flag cover, $\det\Omega$ is isomorphic to $\mathcal L_{\mu_0}$.
\formallabel{GlobalTheory.FlagHodge.lbIso_detOm_Llam}
\end{proposition}

\begin{proof}
Two line bundles with cocycles $(g_{kk'})$ and $(g'_{kk'})$ on the same cover are isomorphic as soon as
there are units $d_k\in\Gamma(U_k,\OO)^\times$ with $g'_{kk'}d_{k'}=g_{kk'}d_k$ on $U_k\cap U_{k'}$; this is
the usual coboundary computation with the cocycle relations. Let $d_k$ be the determinant of the $n\times n$
matrix whose columns are the columns of all $u_k(\tau)$; its value at a point is the determinant of the
flag vectors in the basis of chart $k$.

\emph{Eigenframes.} For every point $z$ of $U_k$ with values in a ring $S$ there is a homomorphism
$\phi\colon F\to S$ such that, for every $\tau$, the flag vectors $v_k(\tau,i)(z)$ lie in the eigenspace of
the Lie algebra at $\pi(z)$ on which $\OO_E$ acts through $\phi\circ\tau$, and form a basis of it
(\lref{GlobalTheory.FlagHodge.exists_eigFrame}, which deduces this from the corresponding property over a
dense open subset of $\mathrm{Fl}$ and the normality of $\mathrm{Fl}$). The Lie algebra is the direct sum of
these eigenspaces (\lref{GlobalTheory.T.existsUnique_embedding_eigenspace_sum}), so all flag vectors
together form a basis (\lref{GlobalTheory.ShFam.exists_eigBasis}); hence $d_k(z)$, the determinant of a
basis in another basis, is a unit, and $d_k$ is a unit on $U_k$ (a function whose values at all points are
units is a unit).

\emph{The coboundary relation.} Let $z$ be a point of $U_k\cap U_{k'}$ and $f$, $f'$ the eigenframes given by
the flag vectors of $k$ and $k'$. By the displayed relation and the linear independence of the flag
vectors, $f'(\tau)=f(\tau)B_{kk'}(\tau)(z)$; so the matrix of $f'$ in the basis $f$ is block diagonal with the
upper triangular blocks $B_{kk'}(\tau)(z)$, whose determinant is $\prod_{\tau,i}T_{kk'}(\tau,i)(z)$. By the
multiplicativity of determinants,
\[
\det\mathrm{chg}_{kk'}(z)\,d_{k'}(z)=\det\nolimits_{M_k}(f')=\det\nolimits_{M_k}(f)\prod_{\tau,i}T_{kk'}(\tau,i)(z)
=d_k(z)\prod_{\tau,i}T_{kk'}(\tau,i)(z),
\]
where $\det_{M_k}$ is the determinant in the basis of chart $k$
(\lref{GlobalTheory.ShFam.ev_detChg_mul_frameDet}, \lref{GlobalTheory.ShFam.frameDet_change}). A
function is determined by its values at all points with values in rings, so
$\bigl(\prod_{\tau,i}T_{kk'}(\tau,i)\bigr)^{-1}d_{k'}=(\det\mathrm{chg}_{kk'})^{-1}d_k$, which is the
coboundary relation between the cocycles of $\mathcal L_{\mu_0}$ and of $\det\Omega$.
\end{proof}

\subsubsection*{Flat line generators}
Fix a CM host datum with field $E$ and $d$ (so that $\sqrt{-d}\in E$), a unitary datum of signature
$(p,q)$ with $p+q\ge1$, a level, the moduli space $X^o$ and a coefficient datum with field $F$, and an
embedding $\tau\colon E\to\C$ which maps $\sqrt{-d}$ to the chosen root of $-d$ in $\C$.

\begin{definition}[Frame functions]\label{pf:covol:def:framefn}
Let $U\subseteq X^o_F$.
\begin{enumerate}[label=(\alph*)]
\item A \emph{frame function} $\varphi$ assigns to every field $F'$, every $F'$-point $y$ of $X^o_F$ and all
  families $v=(v_1,\dots,v_p)$, $w=(w_1,\dots,w_q)$ of elements of the Lie algebra of the universal abelian
  scheme at $y$ an element $\varphi(y,v,w)\in F'$. The triple $(y,v,w)$ is \emph{framed over $U$} (for
  $\beta\colon F\to F'$) if $y$ lies over $\beta$, its image lies in $U$, and $v$ and $w$ are bases of the
  eigenspaces of the Lie algebra on which $\OO_E$ acts through $\beta\circ\tau$ and
  $\beta\circ\bar\tau$ respectively.
\item $\varphi$ is a \emph{line function on $U$} if it does not vanish at framed triples, satisfies
  $\varphi(y,v',w')\det A=\varphi(y,v,w)\det B$ when $v'=Av$ and $w'=Bw$, and is compatible with extensions of
  fields.
\item Let $Y$ be a model of $X^o$ over a Noetherian subring $O\subseteq F$ with fraction field $F$, and let
  $Z$ be a component over $Y$ for $U$: a point of $X^o_F$ with values in a field $\Omega$ which maps to the
  generic point of a component of $Y$ and whose closure is $U$, together with a model $V$ over $\Omega$ of the
  Lie algebra of the semi-abelian scheme of $Y$ at the image point, with its $\OO_E$-action; $Z$ is \emph{geometric} if $F$ is algebraically closed in $\Omega$.
  The \emph{eigen-lattices} of $Z$ are the lattices in the eigenspaces of $V$ for $\tau$ and $\bar\tau$
  given by the Lie algebra at the points of the generic fibre. A line function $\varphi$ on $U$ is
  \emph{flat on the model} if there are $a\in\bigwedge^pV$ and $b\in\bigwedge^qV$ which generate the top
  exterior powers of the two eigen-lattices at all points of the generic fibre, such that $\varphi$
  represents $(a,b)$: whenever $a=\alpha\,v_1\wedge\dots\wedge v_p$ and $b=\kappa\,w_1\wedge\dots\wedge w_q$
  with $v,w$ read at $y$, $\varphi(y,v,w)$ is the value of $\alpha/\kappa$ at $y$.
\item For $s\colon F\to\C$ and $t\in\R$, $\varphi$ has \emph{archimedean norm $t$ at $s$} if
  $\log|\varphi(y,v,w)|+\frac12(\log|\det\Gamma|-\log|\det\Delta|)=t$ for all complex triples $(y,v,w)$
  framed over $U$ for $s$, where $\Gamma$ and $\Delta$ are the Gram matrices of $v$ and $w$ for the Hodge
  metric. For a prime $\ell$, $\kappa\colon F\to\C_\ell$ and $t\in\R$, $\varphi$ has \emph{finite norm $t$ at
  $\kappa$} if $\log\|j(\varphi(y,v,w))\|=t$ for every discrete valuation ring $V'$ with fraction field $K'$,
  every embedding $j\colon K'\to\C_\ell$ with $V'=\{\|j(\cdot)\|\le1\}$ and $j\beta=\kappa$, and every triple
  framed over $U$ for $\beta$ at a point $y$ which extends to a semi-abelian scheme over $V'$ such that $v,w$
  are integral eigenframes of it.
\item Let $\mathrm{FU}$ be a family of complex uniformisations of $X^o$ for a type function of signature
  $(p,q)$, with period lattices $\Lambda_i$, and let $i$ be an index whose period points lie in $U$. A line
  function $\varphi$ is a \emph{flat generator} for $\mathrm{FU}$ at $i$ if there are a family $e$ of $p+q$
  linearly independent elements of $\Lambda_i$ over $\OO_E$ and $\rho\in\C$ such that
  $\varphi(y,v,w)\det\binom{C}{F_m}=\rho$ whenever $y$ is the point of a period map $P$, $(y,v,w)$ is framed
  over $U$, $C$ expresses the derivatives of the uniformisation in the directions $v$ through the
  $\tau$-parts of the vectors $P(e_m)$, and $F_m$ expresses those in the directions $w$ through the Hermitian
  pairings $\tau\langle e_m,e_n\rangle$.
\end{enumerate}
\formal{GlobalTheory.CompanionTheory.Ar.Fl.FrameFn}\formal{GlobalTheory.CompanionTheory.Ar.Fl.IsFramed}%
\formal{GlobalTheory.CompanionTheory.Ar.Fl.IsLineFn}\formal{GlobalTheory.CompanionTheory.Ar.Flat.FlatOnModel}%
\formal{GlobalTheory.CompanionTheory.Ar.Fl.HasInfNorm}\formal{GlobalTheory.CompanionTheory.Ar.Fl.HasFinNorm}%
\formal{GlobalTheory.CompanionTheory.Ar.Fl.IsFlatGen}\formal{GlobalTheory.CompanionTheory.Ar.Fl.IsConstMul}
\end{definition}

\emph{Resolution in projective coordinates} is the statement that a closed subscheme $X$ of a projective space
over a field of characteristic $0$, with an open subset as in the statement, has a resolution of
singularities $X'\to X$ by a closed subscheme $X'$ of a product of two projective spaces compatible with the
projection (a form of resolution of singularities, compare \cref{thm:covolume:resolution});
\emph{GAGA for sections} is the conclusion of \cref{pf:covol:thm:gaga}, which is a hypothesis of the next
two theorems.

\begin{theorem}[Flat line generators]\label{pf:covol:thm:flatzero}
Assume resolution in projective coordinates and GAGA for sections. Let $Y$ be a model of $X^o$ over a
Noetherian subring $O\subseteq F$ with fraction field $F$, $U\subseteq X^o_F$ and $Z$ a geometric component
over $Y$ for $U$. Then there is a line function $\varphi$ on $U$ which is flat on the model, and for every
other such $\varphi'$ there is $x\in F^\times$ with $\varphi'(y,v,w)=\beta(x)\varphi(y,v,w)$ for all triples
framed over $U$ for $\beta$.
\formallabel{GlobalTheory.T.flat_zero_from_resolution_and_gaga}
\end{theorem}

\begin{proof}
Let $\sigma\colon F\to\C$ be the given embedding of the coefficient field. Choose a type function of
signature $(p,q)$ (\lref{GlobalTheory.CompanionTheory.Ar.Flat.A0.exists_typeFn}) and a family of complex
uniformisations of $X^o$ at $\sigma$ for it (\lref{GlobalTheory.familyUnif_of_laws}, which constructs it
from the laws of abelian schemes). There is an \emph{analytic unit function} $\Phi$ for $(\tau,\sigma)$: a
function of complex framed triples which does not vanish, transforms like a line function, satisfies
$|\Phi(y,v,w)|^2|\det\Gamma|=|\det\Delta|$ with the Hodge Gram matrices, and is holomorphic in holomorphic
families (\lref{GlobalTheory.T.exists_unit_analytic_function}). The eigen-lattices of $Z$ are vector
bundles of ranks $p$ and $q$ on the generic fibre, which is not empty
(\lref{GlobalTheory.T.eigen_lattice_vector_bundle}); let $b_1$, $b_2$ be bases over $\Omega$ of the two
eigenspaces of $V$. By resolution and GAGA, there is $f\in\Omega$ such that $a=f\,b_{1,1}\wedge\dots\wedge
b_{1,p}$ and $b=b_{2,1}\wedge\dots\wedge b_{2,q}$ generate the top exterior powers of the eigen-lattices at
all points of the generic fibre (\lref{GlobalTheory.T.eigen_frames_analytic}, which in addition relates
$f$ to $\Phi$ at the complex points; only the generation statement is used here). There is a line function
$\varphi$ on $U$ representing $(a,b)$ (\lref{GlobalTheory.T.wedge_eigen_line_generator}); it is flat on
the model by definition. Finally, two line functions which are flat on the model differ by a nonzero
constant of $F$ (\lref{GlobalTheory.T.frame_functions_flat_agree}).
\end{proof}

\begin{theorem}[Norms and flatness of flat line generators]\label{pf:covol:thm:flatfirst}
Assume resolution in projective coordinates and GAGA for sections. In the situation of
\cref{pf:covol:thm:flatzero}, let $\varphi$ be a line function on $U$ which is flat on the model. Then:
\begin{enumerate}[label=(\alph*)]
\item for every $s\colon F\to\C$ there is $t\in\R$ such that $\varphi$ has archimedean norm $t$ at $s$ and at
  $\bar s$;
\item for every family of complex uniformisations $\mathrm{FU}$ for a type function of signature $(p,q)$ and
  every index $i$ whose period points lie in $U$, $\varphi$ is a flat generator for $\mathrm{FU}$ at $i$.
\end{enumerate}
\formallabel{GlobalTheory.T.flat_first_from_resolution_and_gaga}
\end{theorem}

\begin{proof}
(a) Let $s\colon F\to\C$. Then $s(\tau(\sqrt{-d}))$ is one of the two roots of $-d$, as its square is $-d$. Suppose first that it is the
chosen root. Choose a type function and a family of complex uniformisations at $s$ as in the proof of
\cref{pf:covol:thm:flatzero}, and an analytic unit function $\Phi$ for $(\tau,s)$
(\lref{GlobalTheory.T.exists_unit_analytic_function}). As $\varphi$ is flat on the model, there is
$\kappa\ne0$ with $\varphi=\kappa\Phi$ at all complex triples framed over $U$ for $s$
(\lref{GlobalTheory.T.frame_function_analytic_norm}, which uses resolution and GAGA). The Hodge Gram
matrices of frames are nonsingular
(\lref{GlobalTheory.CompanionTheory.Ar.Flat.FlatFirstTheorem.gram_det_ne_zero}), so taking logarithms in
$|\Phi(y,v,w)|^2|\det\Gamma|=|\det\Delta|$ gives
$\log|\varphi(y,v,w)|+\frac12(\log|\det\Gamma|-\log|\det\Delta|)=\log|\kappa|$: $\varphi$ has archimedean norm
$\log|\kappa|$ at $s$. Otherwise $\bar s$ maps $\tau(\sqrt{-d})$ to the chosen root, as complex conjugation
maps it to its negative, so $\varphi$ has an archimedean norm $t$ at $\bar s$. In both cases the archimedean
norm at the conjugate embedding is the same: if a line function has archimedean norm $t$ at $s$, it has
archimedean norm $t$ at $\bar s$ (\lref{GlobalTheory.T.inf_norm_conjugate_invariant}).

(b) The period points lie in $U$ by hypothesis and $\varphi$ is a line function. Let $\Phi$ be an analytic
unit function for $(\tau,\sigma)$ which is flat for $\mathrm{FU}$ at $i$: there are linearly independent
$e_m\in\Lambda_i$ and $\rho\ne0$ with $\Phi(y,v,w)\det\binom{C}{F_m}=\rho$ in the situation of
\cref{pf:covol:def:framefn}(e) (\lref{GlobalTheory.T.exists_unit_analytic_function}). As $\varphi$ is flat
on the model, $\varphi=\kappa\Phi$ at all complex triples framed over $U$ for $\sigma$, for some $\kappa\ne0$
(\lref{GlobalTheory.T.frame_function_analytic_norm}, which uses resolution and GAGA). Hence
$\varphi(y,v,w)\det\binom{C}{F_m}=\kappa\rho$ in that situation, and $\varphi$ is a flat generator with the
basis $e$ and the constant $\kappa\rho$.
\end{proof}

\subsubsection*{The flat degree}
\begin{definition}[Log norms]\label{pf:covol:def:lognorms}
A family of \emph{log norms} $N$ on a number field $k$ consists of real numbers $N_\infty(s)$ for the
embeddings $s\colon k\to\C$ and $N_\ell(j)$ for the primes $\ell$ and the embeddings $j\colon k\to\C_\ell$. It
is \emph{adelic} if $N_\ell=0$ for almost all $\ell$, and its \emph{degree} is
$\deg N=-\frac1{[k:\Q]}\bigl(\sum_sN_\infty(s)+\sum_\ell\sum_jN_\ell(j)\bigr)$ (with the finite sum taken over the
finitely many nonzero terms). For a frame function $\varphi$ with coefficient field $F$, $N$ is the family of
\emph{norms of $\varphi$} if $\varphi$ has archimedean norm $N_\infty(s)$ at every $s$ and finite norm
$N_\ell(\kappa)+q\bigl(\log\|\mathfrak d_E\|_{\kappa\tau}-\log\|d\|_\ell\bigr)$ at every $\kappa\colon
F\to\C_\ell$, where $\mathfrak d_E$ is the different of $E$ and $\log\|\mathfrak d_E\|_{j}$ is the logarithm
of the absolute value of a local generator at the place defined by $j\colon E\to\C_\ell$.
\formal{GlobalTheory.CompanionTheory.Ar.Fl.LogNorms}\formal{GlobalTheory.CompanionTheory.Ar.Fl.HasNorms}
\end{definition}

Let $\chi$ be the Kronecker character of conductor $d$ and
$\log\Omega(E,\chi)=\gamma+\log(2\pi/\mathrm{rd}_E)-\frac{L'}{L}(1,\chi)$ (compare \cref{rem:covolume:omega}),
where $\mathrm{rd}_E=|d_E|^{1/[E:\Q]}$ is the root discriminant.

\begin{lemma}[An adelic computation]\label{pf:covol:lem:adelic}
Let $k$ be a number field, $\tau\colon E\to k$, $N$ log norms on $k$, $\rho\in k^\times$, $a\in E^\times$,
$r\in E$ with $r^2=-d$, $x\in\R$, positive functions $n,n'$ on the embeddings of $k$ with
$-\frac1{[k:\Q]}\sum_s\log n(s)=x=-\frac1{[k:\Q]}\sum_s\log n'(s)$, and $t_\ell,t'_\ell\in E$ for every prime
$\ell$. Write $L_j=\log\|\mathfrak d_E\|_j$ for $j\colon E\to\C_\ell$. Suppose that
\begin{enumerate}[label=(\roman*)]
\item $N_\infty(s)=\log|s\rho|+\frac p2\bigl(\log n(s)+\frac{\log d}2\bigr)-\frac q2\bigl(\log n'(s)+
  \frac{\log d}2\bigr)+\frac12\log|s\tau(a)|$ for all $s$;
\item $\sum_j\log\|j(t'_\ell)\|=\sum_j\log\|j(t_\ell)\|$ and
  $\log\|j(t_\ell)\|+\log\|j(t'_\ell)\|+\log\|j(a)\|=-(p+q)L_j$ for all $\ell$ and $j\colon E\to\C_\ell$;
\item $t_\ell$ is a unit at all places above $\ell$ for almost all $\ell$;
\item $N_\ell(j)+q(L_{j\tau}-\log\|d\|_\ell)=\log\|j\rho\|-\log\|j\tau(t_\ell)\|-qL_{j\tau}+(q-p)
  \bigl(L_{j\tau}-\log\|j\tau(r)\|\bigr)$ for all $\ell$ and $j\colon k\to\C_\ell$.
\end{enumerate}
Then $N$ is adelic and
$\deg N=\frac{p-q}2\bigl(x-\log\mathrm{rd}_E+\log\sqrt d\bigr)+q\bigl(2\log\sqrt d-\log\mathrm{rd}_E\bigr)$.
\formallabel{GlobalTheory.T.hermitian_metric_log_norm_formula}
\end{lemma}

\begin{proof}
By (iv), $N_\ell(j)=\log\|j\rho\|-\log\|j\tau(t_\ell)\|-(p+q)L_{j\tau}+(p-q)\log\|j\tau(r)\|+q\log\|d\|_\ell$.
For almost all $\ell$, the elements $\rho$, $t_\ell$ (by (iii)), $r$ and $d$ are units above $\ell$ and $\ell$
does not divide $d_E$; then $N_\ell=0$, so $N$ is adelic. Let $n_k=[k:\Q]$, $n_E=[E:\Q]$. Every embedding of
$E$ into $\C$ or $\C_\ell$ is $j\tau$ for exactly $n_k/n_E$ embeddings $j$ of $k$. We use the product formula
$\sum_s\log|s\alpha|+\sum_\ell\sum_j\log\|j\alpha\|=0$ for $\alpha\in k^\times$ (and its analogue for $E$ and
for $\Q$), and $\sum_\ell\sum_{j\colon E\to\C_\ell}L_j=-\log|d_E|=-n_E\log\mathrm{rd}_E$. By (ii),
$2\sum_j\log\|j(t_\ell)\|=-\sum_j\log\|j(a)\|-(p+q)\sum_jL_j$, so
\[
\sum_\ell\sum_{j\colon k\to\C_\ell}\log\|j\tau(t_\ell)\|=\frac{n_k}{2n_E}\Bigl(\sum_{\sigma\colon E\to\C}\log|\sigma a|
+(p+q)\,n_E\log\mathrm{rd}_E\Bigr).
\]
Moreover $\sum_\ell\sum_j\log\|j\tau(r)\|=-n_k\log\sqrt d$ and $\sum_\ell\sum_j\log\|d\|_\ell=-n_k\log d$.
Summing, the finite part is
\begin{multline*}
\sum_\ell\sum_jN_\ell(j)=-\sum_s\log|s\rho|-\frac{n_k}{2n_E}\sum_\sigma\log|\sigma a|-\frac{p+q}2n_k\log\mathrm{rd}_E\\
+(p+q)n_k\log\mathrm{rd}_E-(p-q)n_k\log\sqrt d-qn_k\log d .
\end{multline*}
By (i) and the hypothesis on $n,n'$, the archimedean part is
\[
\sum_sN_\infty(s)=\sum_s\log|s\rho|-\frac{p-q}2n_kx+\frac{p-q}2n_k\log\sqrt d+\frac{n_k}{2n_E}\sum_\sigma\log|\sigma a|.
\]
Adding, dividing by $-n_k$ and using $\log d=2\log\sqrt d$,
\[
\deg N=\frac{p-q}2x+\frac{p-q}2\log\sqrt d-\frac{p+q}2\log\mathrm{rd}_E+2q\log\sqrt d,
\]
which is the asserted value.
\end{proof}

\begin{theorem}[The flat degree]\label{pf:covol:thm:flatdeg}
Let $F$ be a number field, $\mathrm{FU}$ a family of complex uniformisations of $X^o$ for a type function of
signature $(p,q)$, $i$ an index, $U\subseteq X^o_F$, and $\varphi$ a flat generator for $\mathrm{FU}$ at $i$
whose family of norms is $N$. Then $N$ is adelic and
\[
\deg N=\frac{p-q}2\,\Re\log\Omega(E,\chi)+q\bigl(2\log\sqrt d-\log\mathrm{rd}_E\bigr).
\]
\formallabel{GlobalTheory.T.FlatDegree_holds}
\end{theorem}

\begin{proof}
Let $e$ and $\rho$ be the basis and the constant of the flat generator, and $a\in E^\times$ the determinant of
the Gram matrix of $e$ for the Hermitian form of the lattice; $a\ne0$, and $e$ is a basis of the lattice
locally at all but finitely many primes (\lref{GlobalTheory.T.gram_det_ne_zero_local_basis_cofinite}).
\begin{enumerate}[label=(\arabic*)]
\item \emph{Chowla--Selberg.} For each of the two square roots $r_\pm$ of $-d$ in $\C$ there are a number
  field $M_\pm\subseteq\C$ and an elliptic curve $Z_\pm$ over $\OO_{M_\pm}$ with complex multiplication by
  $\frac{1+r_\pm}2$, with a complex uniformisation whose differential is $c_\pm z_\pm$ times the invariant
  vector for nonzero $c_\pm,z_\pm$, and with \emph{period} $x$: the average over the complex embeddings of
  $M_\pm$ of minus the logarithm of the squared norm $n_\pm$ of the invariant vector is $x$, where
  $x=\gamma+\log(2\pi/\sqrt d)-\frac{L'}{L}(1,\chi)$ (\lref{GlobalTheory.T.chowla_selberg_period}); $x$ is
  real because $\chi$ is real.
\item \emph{The CM value.} There are a number field $k\subseteq\C$ containing $F$, $M_+$ and $M_-$, and
  $r\in k^\times$ with $r z_+^p=\rho z_-^q$ (\lref{GlobalTheory.T.frame_function_cm_value}).
\item \emph{Archimedean norms.} For every $s\colon k\to\C$ the archimedean norm of $\varphi$ at $s|_F$ is
  $\log|s r|+\frac p2(\log n_+(s)+\frac{\log d}2)-\frac q2(\log n_-(s)+\frac{\log d}2)+\frac12\log|s\tau(a)|$
  (\lref{GlobalTheory.T.frame_function_cm_number_field}); and the average of $-\log n_\pm(s)$ over the
  embeddings $s$ of $k$ equals that over the embeddings of $M_\pm$, which is $x$.
\item \emph{Finite norms.} For every prime $\ell$ there is $t_\ell\in E$ with
  $\log\|j(t_\ell)\|+\log\|j(\bar t_\ell)\|+\log\|j(a)\|=-(p+q)L_j$ for all $j\colon E\to\C_\ell$, with $t_\ell$
  a unit above $\ell$ whenever $e$ is a local basis at $\ell$, and such that the finite norm of $\varphi$ at
  $j|_F$, for $j\colon k\to\C_\ell$, is
  $\log\|jr\|-\log\|j\tau(t_\ell)\|-qL_{j\tau}+(q-p)(L_{j\tau}-\log\|j\tau(\sqrt{-d})\|)$
  (\lref{GlobalTheory.T.frame_function_cm_height}).
\end{enumerate}
Let $N_k$ be the extension of $N$ to $k$, $N_{k,\infty}(s)=N_\infty(s|_F)$ and $N_{k,\ell}(j)=N_\ell(j|_F)$; it has
the same degree as $N$ and is adelic if and only if $N$ is
(each embedding of $F$ extends to $[k:F]$ embeddings of $k$). Since $\sum_j\log\|j(\bar t)\|=\sum_j\log\|j(t)\|$ for $t\in E$ (complex
conjugation permutes the embeddings), (1)--(4) and the definition of the norms of $\varphi$ give the
hypotheses of \cref{pf:covol:lem:adelic} for $N_k$, $\rho=r$, $r=\tau^{-1}$ of the chosen root,
$n=n_+$, $n'=n_-$, $t'_\ell=\bar t_\ell$. Hence $N$ is adelic and $\deg N=\frac{p-q}2(x-\log\mathrm{rd}_E+
\log\sqrt d)+q(2\log\sqrt d-\log\mathrm{rd}_E)$. Finally $x-\log\mathrm{rd}_E+\log\sqrt d=\gamma+
\log(2\pi/\mathrm{rd}_E)-\frac{L'}{L}(1,\chi)=\Re\log\Omega(E,\chi)$
(\lref{GlobalTheory.CompanionTheory.Ar.Fl.flatDegree_logOm_re}).
\end{proof}
